\section{2021年高考题}
\subsection{2021年全国I卷（理）}\label{gk:2021年全国I卷（理）}

\begin{description}
    \item[一、选择题] 本题共12小题，每小题5分，共60分．在每小题给出的四个选项中，只有一项是符合题目要求的．
    \begin{enumerate}[leftmargin=5pt]
        \item 设集合 $M=\{x\mid0<x<4\}$，$N=\bigl\{x\mid\dfrac{1}{3}\leqslant x\leqslant5\bigr\}$，则 $M\cap N=$
       
        {\centering\begin{tabular*}{\linewidth}{@{}@{\extracolsep{\fill}}llll@{}}
          \(\text{A. }\bigl\{x\mid0<x\leqslant\dfrac{1}{3}\bigr\}\)  & \(\text{B. }\bigl\{x\mid\dfrac{1}{3}\leqslant x<4\bigr\}\)  & \(\text{C. }\{x\mid4\leqslant x<5\}\)  & \(\text{D. }\{x\mid0<x\leqslant5\}\)
        \end{tabular*}\par}

        \item 为了解某地农村经济情况，对该地农户家庭年收入进行抽样调查，将农户家庭年收入的调查数据整理得到如下频率分布直方图：
        
        \textsf{[插入图片：频率分布直方图]}

        根据此频率分布直方图，下面结论中不正确的是
       
        {\centering\begin{tabular*}{\linewidth}{@{}@{\extracolsep{\fill}}llll@{}}
          \(\text{A. }\)该地农户家庭年收入低于4.5万元的农户比率估计为$6\%$ \\
          \(\text{B. }\)该地农户家庭年收入不低于10.5万元的农户比率估计为$10\%$ \\
          \(\text{C. }\)估计该地农户家庭年收入的平均值不超过6.5万元 \\
          \(\text{D. }\)估计该地有一半以上的农户，其家庭年收入介于4.5万元至8.5万元之间
        \end{tabular*}\par}

        \item 已知 $(1-\i)^2z=3+2\i$，则 $z=$
       
        {\centering\begin{tabular*}{\linewidth}{@{}@{\extracolsep{\fill}}llll@{}}
          \(\text{A. }-1-\dfrac{3}{2}\i\)  & \(\text{B. }-1+\dfrac{3}{2}\i\)  & \(\text{C. }-\dfrac{3}{2}+\i\)  & \(\text{D. }-\dfrac{3}{2}-\i\)
        \end{tabular*}\par}

        \item 青少年视力是社会普遍关注的问题，视力情况可借助视力表测量．通常用五分记录法和小数记录法记录视力数据，五分记录法的数据 $L$ 和小数记录法的数据 $V$ 满足 $L=5+\lg V$．已知某同学视力的五分记录法的数据为 $4.9$，则其视力的小数记录法的数据约为 $(\sqrt[10]{10}\approx1.259)$
       
        {\centering\begin{tabular*}{\linewidth}{@{}@{\extracolsep{\fill}}llll@{}}
          \(\text{A. }1.5\)  & \(\text{B. }1.2\)  & \(\text{C. }0.8\)  & \(\text{D. }0.6\)
        \end{tabular*}\par}

        \item 已知 $F_1,F_2$ 是双曲线 $C$ 的两个焦点，$P$ 为 $C$ 上一点，且 $\angle F_1PF_2=60^\circ$，$|PF_1|=3|PF_2|$，则 $C$ 的离心率为
       
        {\centering\begin{tabular*}{\linewidth}{@{}@{\extracolsep{\fill}}llll@{}}
          \(\text{A. }\dfrac{\sqrt{7}}{2}\)  & \(\text{B. }\dfrac{\sqrt{13}}{2}\)  & \(\text{C. }\sqrt{7}\)  & \(\text{D. }\sqrt{13}\)
        \end{tabular*}\par}

        \item 在一个正方体中，过顶点 $A$ 的三条棱的中点分别为 $E,F,G$．该正方体截去三棱锥 $A$-$EFG$ 后，所得多面体的三视图中，正视图如右图所示，则相应的侧视图是
        
        \textsf{[插入图片：三视图选项]}

        \item 等比数列 $\{a_n\}$ 的公比为 $q$，前 $n$ 项和为 $S_n$．设甲：$q>0$，乙：$\{S_n\}$ 是递增数列，则
       
        {\centering\begin{tabular*}{\linewidth}{@{}@{\extracolsep{\fill}}llll@{}}
          \(\text{A. }\)甲是乙的充分条件但不是必要条件  & \(\text{B. }\)甲是乙的必要条件但不是充分条件 \\
          \(\text{C. }\)甲是乙的充要条件  & \(\text{D. }\)甲既不是乙的充分条件也不是乙的必要条件
        \end{tabular*}\par}

        \item 2020年12月8日，中国和尼泊尔联合公布珠穆朗玛峰最新高程为 $8848.86$（单位：m），三角高程测量法是珠峰高程测量方法之一．如图是三角高程测量法的一个示意图，现有 $A,B,C$ 三点，且 $A,B,C$ 在同一水平面上的投影 $A',B',C'$ 满足 $\angle A'C'B'=45^\circ$，$\angle A'B'C'=60^\circ$．由 $C$ 点测得 $B$ 点的仰角为 $15^\circ$，$BB'$ 与 $CC'$ 的差为 $100$；由 $B$ 点测得 $A$ 点的仰角为 $45^\circ$，则 $A,C$ 两点到水平面 $A'B'C'$ 的高度差 $AA'-CC'$ 约为 $(\sqrt{3}\approx1.732)$
        
        \textsf{[插入图片：三角高程测量示意图]}

        {\centering\begin{tabular*}{\linewidth}{@{}@{\extracolsep{\fill}}llll@{}}
          \(\text{A. }346\)  & \(\text{B. }373\)  & \(\text{C. }446\)  & \(\text{D. }473\)
        \end{tabular*}\par}

        \item 若 $\alpha\in\bigl(0,\dfrac{\pi}{2}\bigr)$，$\tan2\alpha=\dfrac{\cos\alpha}{2-\sin\alpha}$，则 $\tan\alpha=$
       
        {\centering\begin{tabular*}{\linewidth}{@{}@{\extracolsep{\fill}}llll@{}}
          \(\text{A. }\dfrac{\sqrt{15}}{15}\)  & \(\text{B. }\dfrac{\sqrt{5}}{5}\)  & \(\text{C. }\dfrac{\sqrt{5}}{3}\)  & \(\text{D. }\dfrac{\sqrt{15}}{3}\)
        \end{tabular*}\par}

        \item 将4个1和2个0随机排成一行，则2个0不相邻的概率为
       
        {\centering\begin{tabular*}{\linewidth}{@{}@{\extracolsep{\fill}}llll@{}}
          \(\text{A. }\dfrac{1}{3}\)  & \(\text{B. }\dfrac{2}{5}\)  & \(\text{C. }\dfrac{2}{3}\)  & \(\text{D. }\dfrac{4}{5}\)
        \end{tabular*}\par}

        \item 已知 $A,B,C$ 是半径为 $1$ 的球 $O$ 的球面上的三个点，且 $AC\perp BC$，$AC=BC=1$，则三棱锥 $O$-$ABC$ 的体积为
       
        {\centering\begin{tabular*}{\linewidth}{@{}@{\extracolsep{\fill}}llll@{}}
          \(\text{A. }\dfrac{\sqrt{2}}{12}\)  & \(\text{B. }\dfrac{\sqrt{3}}{12}\)  & \(\text{C. }\dfrac{\sqrt{2}}{4}\)  & \(\text{D. }\dfrac{\sqrt{3}}{4}\)
        \end{tabular*}\par}

        \item 设函数 $f(x)$ 的定义域为 $\mathbb{R}$，$f(x+1)$ 为奇函数，$f(x+2)$ 为偶函数，当 $x\in[1,2]$ 时，$f(x)=ax^2+b$．若 $f(0)+f(3)=6$，则 $f\bigl(\dfrac{9}{2}\bigr)=$
       
        {\centering\begin{tabular*}{\linewidth}{@{}@{\extracolsep{\fill}}llll@{}}
          \(\text{A. }-\dfrac{9}{4}\)  & \(\text{B. }-\dfrac{3}{2}\)  & \(\text{C. }\dfrac{7}{4}\)  & \(\text{D. }\dfrac{5}{2}\)
        \end{tabular*}\par}
    \end{enumerate}
\end{description}

\begin{description}
    \item[二、填空题] 本题共4小题，每小题5分，共20分．
    \begin{enumerate}[leftmargin=5pt]
      \addtocounter{enumi}{+12}
        \item 曲线 $y=\dfrac{2x-1}{x+2}$ 在点 $(-1,-3)$ 处的切线方程为 \rule[-0.3ex]{3em}{0.4pt}．
        \item 已知向量 $a=(1,3)$，$b=(3,4)$，若 $(a-\lambda b)\perp b$，则 $\lambda=$ \rule[-0.3ex]{3em}{0.4pt}．
        \item 已知 $F_1,F_2$ 为椭圆 $C:\dfrac{x^2}{16}+\dfrac{y^2}{4}=1$ 的两个焦点，$P,Q$ 为 $C$ 上关于坐标原点对称的两点，且 $|PQ|=|F_1F_2|$，则四边形 $PF_1QF_2$ 的面积为 \rule[-0.3ex]{3em}{0.4pt}．
        \item 已知函数 $f(x)=2\cos(\omega x+\varphi)$ 的部分图像如图所示，则满足条件 $\bigl(f(x)-f\bigl(-\dfrac{7\pi}{4}\bigr)\bigr)\bigl(f(x)-f\bigl(\dfrac{4\pi}{3}\bigr)\bigr)>0$ 的最小正整数 $x$ 为 \rule[-0.3ex]{3em}{0.4pt}．
        
        \textsf{[插入图片：函数图像]}
    \end{enumerate}
\end{description}

\begin{description}
    \item[三、解答题] 本题共7小题，共70分．解答应写出文字说明、证明过程或演算步骤．第17\textasciitilde21题为必考题，第22、23题为选考题．
    \begin{enumerate}[leftmargin=5pt]
      \addtocounter{enumi}{+16}
        \item （12分）

        甲、乙两台机床生产同种产品，产品按质量分为一级品和二级品，为了比较两台机床产品的质量，分别用两台机床各生产了200件产品，产品的质量情况统计如下表：

        {\centering\begin{tabular}{c|c|c|c}
        \hline
        & 一级品 & 二级品 & 合计 \\
        \hline
        甲机床 & 150 & 50 & 200 \\
        \hline
        乙机床 & 120 & 80 & 200 \\
        \hline
        合计 & 270 & 130 & 400 \\
        \hline
        \end{tabular}\par}

        \begin{enumerate}[label=(\arabic*)]
            \item 甲机床、乙机床生产的产品中一级品的频率分别是多少？
            \item 能否有 $99\%$ 的把握认为甲机床的产品质量与乙机床的产品质量有差异？

            附：$K^2=\dfrac{n(ad-bc)^2}{(a+b)(c+d)(a+c)(b+d)}$，$P(K^2\geqslant k)$：0.050--3.841，0.010--6.635，0.001--10.828．
        \end{enumerate}

        \item （12分）

        已知数列 $\{a_n\}$ 的各项均为正数，记 $S_n$ 为 $\{a_n\}$ 的前 $n$ 项和，从下面\ding{172}\ding{173}\ding{174}中选取两个作为条件，证明另外一个成立．

        \ding{172}数列 $\{a_n\}$ 是等差数列；\ding{173}数列 $\{\sqrt{S_n}\}$ 是等差数列；\ding{174}$a_2=3a_1$．

        注：若选择不同的组合分别解答，则按第一个解答计分．

        \item （12分）

        已知直三棱柱 $ABC$-$A_1B_1C_1$ 中，侧面 $AA_1B_1B$ 为正方形，$AB=BC=2$，$E,F$ 分别为 $AC$ 和 $CC_1$ 的中点，$D$ 为棱 $A_1B_1$ 上的点，$BF\perp A_1B_1$．

        \textsf{[插入图片：直三棱柱示意图]}

        \begin{enumerate}[label=(\arabic*)]
            \item 证明：$BF\perp DE$；
            \item 当 $B_1D$ 为何值时，面 $BB_1C_1C$ 与面 $DFE$ 所成的二面角的正弦值最小？
        \end{enumerate}

        \item （12分）

        抛物线 $C$ 的顶点为坐标原点 $O$，焦点在 $x$ 轴上，直线 $l:x=1$ 交 $C$ 于 $P,Q$ 两点，且 $OP\perp OQ$．已知点 $M(2,0)$，且 $\odot M$ 与 $l$ 相切．

        \begin{enumerate}[label=(\arabic*)]
            \item 求 $C$，$\odot M$ 的方程；
            \item 设 $A_1,A_2,A_3$ 是 $C$ 上的三个点，直线 $A_1A_2,A_1A_3$ 均与 $\odot M$ 相切．判断直线 $A_2A_3$ 与 $\odot M$ 的位置关系，并说明理由．
        \end{enumerate}

        \item （12分）

        已知 $a>0$ 且 $a\neq1$，函数 $f(x)=\dfrac{x^a}{a^x}\,(x>0)$．

        \begin{enumerate}[label=(\arabic*)]
            \item 当 $a=2$ 时，求 $f(x)$ 的单调区间；
            \item 若曲线 $y=f(x)$ 与直线 $y=1$ 有且仅有两个交点，求 $a$ 的取值范围．
        \end{enumerate}

        \item （10分）[选修4-4：坐标系与参数方程]

        在直角坐标系 $xOy$ 中，$\odot C$ 的圆心为 $C(2,1)$，半径为 $1$．

        \begin{enumerate}[label=(\arabic*)]
            \item 写出 $\odot C$ 的一个参数方程；
            \item 过点 $F(4,1)$ 作 $\odot C$ 的两条切线．以坐标原点为极点，$x$ 轴正半轴为极轴建立极坐标系，求这两条切线的极坐标方程．
        \end{enumerate}

        \item （10分）[选修4-5：不等式选讲]

        已知函数 $f(x)=|x-a|+|x+3|$．

        \begin{enumerate}[label=(\arabic*)]
            \item 当 $a=1$ 时，求不等式 $f(x)\geqslant6$ 的解集；
            \item 若 $f(x)>-a$，求 $a$ 的取值范围．
        \end{enumerate}
    \end{enumerate}
\end{description}
\clearpage

\subsection{2021年全国I卷（文）}\label{gk:2021年全国I卷（文）}

\begin{description}
    \item[一、选择题] 本题共12小题，每小题5分，共60分．在每小题给出的四个选项中，只有一项是符合题目要求的．
    \begin{enumerate}[leftmargin=5pt]
        \item 已知全集 $U=\{1,2,3,4,5\}$，集合 $M=\{1,2\}$，$N=\{3,4\}$，则 $\complement_U(M\cup N)=$
       
        {\centering\begin{tabular*}{\linewidth}{@{}@{\extracolsep{\fill}}llll@{}}
          \(\text{A. }\{5\}\)  & \(\text{B. }\{1,2\}\)  & \(\text{C. }\{3,4\}\)  & \(\text{D. }\{1,2,3,4\}\)
        \end{tabular*}\par}

        \item 设 $\i z=4+3\i$，则 $z=$
       
        {\centering\begin{tabular*}{\linewidth}{@{}@{\extracolsep{\fill}}llll@{}}
          \(\text{A. }-3-4\i\)  & \(\text{B. }-3+4\i\)  & \(\text{C. }3-4\i\)  & \(\text{D. }3+4\i\)
        \end{tabular*}\par}

        \item 已知命题 $p:\exists x\in\mathbb{R},\sin x<1$；命题 $q:\forall x\in\mathbb{R},\e^{|x|}\geqslant1$，则下列命题中为真命题的是
       
        {\centering\begin{tabular*}{\linewidth}{@{}@{\extracolsep{\fill}}llll@{}}
          \(\text{A. }p\land q\)  & \(\text{B. }\neg p\land q\)  & \(\text{C. }p\land\neg q\)  & \(\text{D. }\neg(p\lor q)\)
        \end{tabular*}\par}

        \item 函数 $f(x)=\sin\dfrac{x}{3}+\cos\dfrac{x}{3}$ 的最小正周期和最大值分别是
       
        {\centering\begin{tabular*}{\linewidth}{@{}@{\extracolsep{\fill}}llll@{}}
          \(\text{A. }3\pi$和$\sqrt{2}\)  & \(\text{B. }3\pi$和$2\)  & \(\text{C. }6\pi$和$\sqrt{2}\)  & \(\text{D. }6\pi$和$2\)
        \end{tabular*}\par}

        \item 若 $x,y$ 满足约束条件 $\begin{cases}x+y\geqslant4,\\x-y\leqslant2,\\y\leqslant3,\end{cases}$ 则 $z=3x+y$ 的最小值为
       
        {\centering\begin{tabular*}{\linewidth}{@{}@{\extracolsep{\fill}}llll@{}}
          \(\text{A. }18\)  & \(\text{B. }10\)  & \(\text{C. }6\)  & \(\text{D. }4\)
        \end{tabular*}\par}

        \item $\cos^2\dfrac{\pi}{12}-\cos^2\dfrac{5\pi}{12}=$
       
        {\centering\begin{tabular*}{\linewidth}{@{}@{\extracolsep{\fill}}llll@{}}
          \(\text{A. }\dfrac{1}{2}\)  & \(\text{B. }\dfrac{\sqrt{3}}{3}\)  & \(\text{C. }\dfrac{\sqrt{2}}{2}\)  & \(\text{D. }\dfrac{\sqrt{3}}{2}\)
        \end{tabular*}\par}

        \item 在区间 $(0,\dfrac{1}{2})$ 随机取 $1$ 个数，则取到的数小于 $\dfrac{1}{3}$ 的概率为
       
        {\centering\begin{tabular*}{\linewidth}{@{}@{\extracolsep{\fill}}llll@{}}
          \(\text{A. }\dfrac{3}{4}\)  & \(\text{B. }\dfrac{2}{3}\)  & \(\text{C. }\dfrac{1}{3}\)  & \(\text{D. }\dfrac{1}{6}\)
        \end{tabular*}\par}

        \item 下列函数中最小值为 $4$ 的是
       
        {\centering\begin{tabular*}{\linewidth}{@{}@{\extracolsep{\fill}}llll@{}}
          \(\text{A. }y=x^2+2x+4\)  & \(\text{B. }y=|\sin x|+\dfrac{4}{|\sin x|}\)  & \(\text{C. }y=2^x+2^{2-x}\)  & \(\text{D. }y=\ln x+\dfrac{4}{\ln x}\)
        \end{tabular*}\par}

        \item 设函数 $f(x)=\dfrac{1-x}{1+x}$，则下列函数中为奇函数的是
       
        {\centering\begin{tabular*}{\linewidth}{@{}@{\extracolsep{\fill}}llll@{}}
          \(\text{A. }f(x-1)-1\)  & \(\text{B. }f(x-1)+1\)  & \(\text{C. }f(x+1)-1\)  & \(\text{D. }f(x+1)+1\)
        \end{tabular*}\par}

        \item 在正方体 $ABCD$-$A_1B_1C_1D_1$ 中，$P$ 为 $B_1D_1$ 的中点，则直线 $PB$ 与 $AD_1$ 所成的角为
       
        {\centering\begin{tabular*}{\linewidth}{@{}@{\extracolsep{\fill}}llll@{}}
          \(\text{A. }\dfrac{\pi}{2}\)  & \(\text{B. }\dfrac{\pi}{3}\)  & \(\text{C. }\dfrac{\pi}{4}\)  & \(\text{D. }\dfrac{\pi}{6}\)
        \end{tabular*}\par}

        \item 设 $B$ 是椭圆 $C:\dfrac{x^2}{5}+y^2=1$ 的上顶点，点 $P$ 在 $C$ 上，则 $|PB|$ 的最大值为
       
        {\centering\begin{tabular*}{\linewidth}{@{}@{\extracolsep{\fill}}llll@{}}
          \(\text{A. }\dfrac{5}{2}\)  & \(\text{B. }\sqrt{6}\)  & \(\text{C. }\sqrt{5}\)  & \(\text{D. }2\)
        \end{tabular*}\par}

        \item 设 $a\neq0$，若 $x=a$ 为函数 $f(x)=a(x-a)^2(x-b)$ 的极大值点，则
       
        {\centering\begin{tabular*}{\linewidth}{@{}@{\extracolsep{\fill}}llll@{}}
          \(\text{A. }a<b\)  & \(\text{B. }a>b\)  & \(\text{C. }ab<a^2\)  & \(\text{D. }ab>a^2\)
        \end{tabular*}\par}
    \end{enumerate}
\end{description}

\begin{description}
    \item[二、填空题] 本题共4小题，每小题5分，共20分．
    \begin{enumerate}[leftmargin=5pt]
      \addtocounter{enumi}{+12}
        \item 已知向量 $a=(2,5)$，$b=(\lambda,4)$，若 $a\parallel b$，则 $\lambda=$ \rule[-0.3ex]{3em}{0.4pt}．
        \item 双曲线 $\dfrac{x^2}{4}-\dfrac{y^2}{5}=1$ 的右焦点到直线 $x+2y-8=0$ 的距离为 \rule[-0.3ex]{3em}{0.4pt}．
        \item 记 $\triangle ABC$ 的内角 $A,B,C$ 的对边分别为 $a,b,c$，面积为 $\sqrt{3}$，$B=60^\circ$，$a^2+c^2=3ac$，则 $b=$ \rule[-0.3ex]{3em}{0.4pt}．
        \item 以图\ding{172}为正视图，在图\ding{173}\ding{174}\ding{175}\ding{176}中选两个分别作为侧视图和俯视图，组成某个三棱锥的三视图，则所选侧视图和俯视图的编号依次为 \rule[-0.3ex]{3em}{0.4pt}（写出符合要求的一组答案即可）．
        
        \textsf{[插入图片：三视图选项]}
    \end{enumerate}
\end{description}

\begin{description}
    \item[三、解答题] 本题共7小题，共70分．解答应写出文字说明、证明过程或演算步骤．第17\textasciitilde21题为必考题，第22、23题为选考题．
    \begin{enumerate}[leftmargin=5pt]
      \addtocounter{enumi}{+16}
        \item （12分）

        某厂研制了一种生产高精产品的设备，为检验新设备生产产品的某项指标有无提高，用一台旧设备和一台新设备各生产了10件产品，得到各件产品该项指标数据如下：

        \textsf{[插入表格：新旧设备数据]}

        旧设备和新设备生产产品的该项指标的样本平均数分别记为 $\bar{x}$ 和 $\bar{y}$，样本方差分别记为 $s_1^2$ 和 $s_2^2$．

        \begin{enumerate}[label=(\arabic*)]
            \item 求 $\bar{x},\bar{y},s_1^2,s_2^2$；
            \item 判断新设备生产产品的该项指标的均值较旧设备是否有显著提高（如果 $\bar{y}-\bar{x}\geqslant2\sqrt{\dfrac{s_1^2+s_2^2}{10}}$，则认为新设备生产产品的该项指标的均值较旧设备有显著提高，否则不认为有显著提高）．
        \end{enumerate}

        \item （12分）

        如图，四棱锥 $P$-$ABCD$ 的底面是矩形，$PD\perp$ 底面 $ABCD$，$M$ 为 $BC$ 的中点，且 $PB\perp AM$．

        \textsf{[插入图片：四棱锥示意图]}

        \begin{enumerate}[label=(\arabic*)]
            \item 证明：平面 $PAM\perp$ 平面 $PBD$；
            \item 若 $PD=DC=1$，求四棱锥 $P$-$ABCD$ 的体积．
        \end{enumerate}

        \item （12分）

        设 $\{a_n\}$ 是首项为 $1$ 的等比数列，数列 $\{b_n\}$ 满足 $b_n=\dfrac{na_n}{3}$．已知 $a_1,3a_2,9a_3$ 成等差数列．

        \begin{enumerate}[label=(\arabic*)]
            \item 求 $\{a_n\}$ 和 $\{b_n\}$ 的通项公式；
            \item 记 $S_n$ 和 $T_n$ 分别为 $\{a_n\}$ 和 $\{b_n\}$ 的前 $n$ 项和．证明：$T_n<\dfrac{S_n}{2}$．
        \end{enumerate}

        \item （12分）

        已知抛物线 $C:y^2=2px\,(p>0)$ 的焦点 $F$ 到准线的距离为 $2$．

        \begin{enumerate}[label=(\arabic*)]
            \item 求 $C$ 的方程；
            \item 已知 $O$ 为坐标原点，点 $P$ 在 $C$ 上，点 $Q$ 满足 $\overrightarrow{PQ}=9\overrightarrow{QF}$，求直线 $OQ$ 斜率的最大值．
        \end{enumerate}

        \item （12分）

        已知函数 $f(x)=x^3-x^2+ax+1$．

        \begin{enumerate}[label=(\arabic*)]
            \item 讨论 $f(x)$ 的单调性；
            \item 求曲线 $y=f(x)$ 过坐标原点的切线与曲线 $y=f(x)$ 的公共点的坐标．
        \end{enumerate}

        \item （10分）[选修4-4：坐标系与参数方程]

        在直角坐标系 $xOy$ 中，$\odot C$ 的圆心为 $C(2,1)$，半径为 $1$．

        \begin{enumerate}[label=(\arabic*)]
            \item 写出 $\odot C$ 的一个参数方程；
            \item 过点 $F(4,1)$ 作 $\odot C$ 的两条切线．以坐标原点为极点，$x$ 轴正半轴为极轴建立极坐标系，求这两条切线的极坐标方程．
        \end{enumerate}

        \item （10分）[选修4-5：不等式选讲]

        已知函数 $f(x)=|x-a|+|x+3|$．

        \begin{enumerate}[label=(\arabic*)]
            \item 当 $a=1$ 时，求不等式 $f(x)\geqslant6$ 的解集；
            \item 若 $f(x)>-a$，求 $a$ 的取值范围．
        \end{enumerate}
    \end{enumerate}
\end{description}
\clearpage

\subsection{2021年全国II卷（理）}\label{gk:2021年全国II卷（理）}

\begin{description}
    \item[一、选择题] 本题共12小题，每小题5分，共60分．在每小题给出的四个选项中，只有一项是符合题目要求的．
    \begin{enumerate}[leftmargin=5pt]
        \item 设集合 $M=\{1,2,3,4,5\}$，$N=\{2,4,6\}$，则 $M\cap N=$
       
        {\centering\begin{tabular*}{\linewidth}{@{}@{\extracolsep{\fill}}llll@{}}
          \(\text{A. }\{2,4\}\)  & \(\text{B. }\{1,3,5\}\)  & \(\text{C. }\{1,2,3,4,5,6\}\)  & \(\text{D. }\{6\}\)
        \end{tabular*}\par}

        \item 为了解某地农村经济情况，对该地农户家庭年收入进行抽样调查，将农户家庭年收入的调查数据整理得到频率分布直方图，根据此频率分布直方图，下面结论中不正确的是
        
        \textsf{[插入图片：频率分布直方图]}

        {\centering\begin{tabular*}{\linewidth}{@{}@{\extracolsep{\fill}}llll@{}}
          \(\text{A. }\)低于4.5万元的农户比率估计为$6\%$ \\
          \(\text{B. }\)不低于10.5万元的农户比率估计为$10\%$ \\
          \(\text{C. }\)平均值不超过6.5万元 \\
          \(\text{D. }\)一半以上农户年收入介于4.5万至8.5万
        \end{tabular*}\par}

        \item 已知 $(1-\i)^2z=3+2\i$，则 $z=$
       
        {\centering\begin{tabular*}{\linewidth}{@{}@{\extracolsep{\fill}}llll@{}}
          \(\text{A. }-1-\dfrac{3}{2}\i\)  & \(\text{B. }-1+\dfrac{3}{2}\i\)  & \(\text{C. }-\dfrac{3}{2}+\i\)  & \(\text{D. }-\dfrac{3}{2}-\i\)
        \end{tabular*}\par}

        \item 设 $a=2\ln1.01$，$b=\ln1.02$，$c=\sqrt{1.04}-1$，则
       
        {\centering\begin{tabular*}{\linewidth}{@{}@{\extracolsep{\fill}}llll@{}}
          \(\text{A. }a<b<c\)  & \(\text{B. }b<c<a\)  & \(\text{C. }b<a<c\)  & \(\text{D. }c<a<b\)
        \end{tabular*}\par}

        \item 在正方体 $ABCD$-$A_1B_1C_1D_1$ 中，$P$ 为 $B_1D_1$ 的中点，则直线 $PB$ 与 $AD_1$ 所成的角为
       
        {\centering\begin{tabular*}{\linewidth}{@{}@{\extracolsep{\fill}}llll@{}}
          \(\text{A. }\dfrac{\pi}{2}\)  & \(\text{B. }\dfrac{\pi}{3}\)  & \(\text{C. }\dfrac{\pi}{4}\)  & \(\text{D. }\dfrac{\pi}{6}\)
        \end{tabular*}\par}

        \item 将5名北京冬奥会志愿者分配到花样滑冰、短道速滑、冰球和冰壶4个项目进行培训，每名志愿者只分配到1个项目，每个项目至少分配1名志愿者，则不同的分配方案共有
       
        {\centering\begin{tabular*}{\linewidth}{@{}@{\extracolsep{\fill}}llll@{}}
          \(\text{A. }60\)种  & \(\text{B. }120\)种  & \(\text{C. }240\)种  & \(\text{D. }480\)种
        \end{tabular*}\par}

        \item 把函数 $y=f(x)$ 图象上所有点的横坐标缩短到原来的 $\dfrac{1}{2}$ 倍，纵坐标不变，再把所得曲线向右平移 $\dfrac{\pi}{3}$ 个单位长度，得到函数 $y=\sin\bigl(x-\dfrac{\pi}{4}\bigr)$ 的图像，则 $f(x)=$
       
        {\centering\begin{tabular*}{\linewidth}{@{}@{\extracolsep{\fill}}llll@{}}
          \(\text{A. }\sin\bigl(\dfrac{x}{2}-\dfrac{7\pi}{12}\bigr)\)  & \(\text{B. }\sin\bigl(\dfrac{x}{2}+\dfrac{\pi}{12}\bigr)\)  & \(\text{C. }\sin\bigl(2x-\dfrac{7\pi}{12}\bigr)\)  & \(\text{D. }\sin\bigl(2x+\dfrac{\pi}{12}\bigr)\)
        \end{tabular*}\par}

        \item 在区间 $(0,1)$ 与 $(1,2)$ 中各随机取 $1$ 个数，则两数之和大于 $\dfrac{7}{4}$ 的概率为
       
        {\centering\begin{tabular*}{\linewidth}{@{}@{\extracolsep{\fill}}llll@{}}
          \(\text{A. }\dfrac{7}{9}\)  & \(\text{B. }\dfrac{23}{32}\)  & \(\text{C. }\dfrac{9}{32}\)  & \(\text{D. }\dfrac{2}{9}\)
        \end{tabular*}\par}

        \item 魏晋时期刘徽撰写的《海岛算经》是关于测量的数学著作，其中第一题是测量海岛的高．如图，点 $E,H,G$ 在水平线 $AC$ 上，$DE$ 和 $FG$ 是两个垂直于水平面且等高的测量标杆的高度，称为"表高"，$EG$ 称为"表距"，$GC$ 和 $EH$ 都称为"表目距"，$GC$ 与 $EH$ 的差称为"表目距的差"．则海岛的高 $AB=$
        
        \textsf{[插入图片：海岛算经示意图]}

        {\centering\begin{tabular*}{\linewidth}{@{}@{\extracolsep{\fill}}llll@{}}
          \(\text{A. }\dfrac{\text{表高}\times\text{表距}}{\text{表目距的差}}+\text{表高}\)  & \(\text{B. }\dfrac{\text{表高}\times\text{表距}}{\text{表目距的差}}-\text{表高}\) \\
          \(\text{C. }\dfrac{\text{表高}\times\text{表距}}{\text{表目距的差}}+\text{表距}\)  & \(\text{D. }\dfrac{\text{表高}\times\text{表距}}{\text{表目距的差}}-\text{表距}\)
        \end{tabular*}\par}

        \item 设 $a\neq0$，若 $x=a$ 为函数 $f(x)=a(x-a)^2(x-b)$ 的极大值点，则
       
        {\centering\begin{tabular*}{\linewidth}{@{}@{\extracolsep{\fill}}llll@{}}
          \(\text{A. }a<b\)  & \(\text{B. }a>b\)  & \(\text{C. }ab<a^2\)  & \(\text{D. }ab>a^2\)
        \end{tabular*}\par}

        \item 设 $B$ 是椭圆 $C:\dfrac{x^2}{a^2}+\dfrac{y^2}{b^2}=1\,(a>b>0)$ 的上顶点，若 $C$ 上的任意一点 $P$ 都满足 $|PB|\leqslant2b$，则 $C$ 的离心率的取值范围是
       
        {\centering\begin{tabular*}{\linewidth}{@{}@{\extracolsep{\fill}}llll@{}}
          \(\text{A. }\bigl[\dfrac{\sqrt{2}}{2},1\bigr)\)  & \(\text{B. }\bigl[\dfrac{1}{2},1\bigr)\)  & \(\text{C. }\bigl(0,\dfrac{\sqrt{2}}{2}\bigr]\)  & \(\text{D. }\bigl(0,\dfrac{1}{2}\bigr]\)
        \end{tabular*}\par}

        \item 设 $a=2\ln1.01$，$b=\ln1.02$，$c=\sqrt{1.04}-1$，则
       
        {\centering\begin{tabular*}{\linewidth}{@{}@{\extracolsep{\fill}}llll@{}}
          \(\text{A. }a<b<c\)  & \(\text{B. }b<c<a\)  & \(\text{C. }b<a<c\)  & \(\text{D. }c<a<b\)
        \end{tabular*}\par}
    \end{enumerate}
\end{description}

\begin{description}
    \item[二、填空题] 本题共4小题，每小题5分，共20分．
    \begin{enumerate}[leftmargin=5pt]
      \addtocounter{enumi}{+12}
        \item 已知双曲线 $C:\dfrac{x^2}{m}-y^2=1\,(m>0)$ 的一条渐近线为 $\sqrt{3}x+my=0$，则 $C$ 的焦距为 \rule[-0.3ex]{3em}{0.4pt}．
        \item 已知向量 $a=(1,3)$，$b=(3,4)$，若 $(a-\lambda b)\perp b$，则 $\lambda=$ \rule[-0.3ex]{3em}{0.4pt}．
        \item 记 $\triangle ABC$ 的内角 $A,B,C$ 的对边分别为 $a,b,c$，面积为 $\sqrt{3}$，$B=60^\circ$，$a^2+c^2=3ac$，则 $b=$ \rule[-0.3ex]{3em}{0.4pt}．
        \item 以图\ding{172}为正视图，在图\ding{173}\ding{174}\ding{175}\ding{176}中选两个分别作为侧视图和俯视图，组成某个三棱锥的三视图，则所选侧视图和俯视图的编号依次为 \rule[-0.3ex]{3em}{0.4pt}．
        
        \textsf{[插入图片：三视图选项]}
    \end{enumerate}
\end{description}

\begin{description}
    \item[三、解答题] 本题共7小题，共70分．解答应写出文字说明、证明过程或演算步骤．第17\textasciitilde21题为必考题，第22、23题为选考题．
    \begin{enumerate}[leftmargin=5pt]
      \addtocounter{enumi}{+16}
        \item （12分）

        记 $S_n$ 为数列 $\{a_n\}$ 的前 $n$ 项和，$b_n$ 为数列 $\{S_n\}$ 的前 $n$ 项积，已知 $\dfrac{2}{S_n}+\dfrac{1}{b_n}=2$．

        \begin{enumerate}[label=(\arabic*)]
            \item 证明：数列 $\{b_n\}$ 是等差数列；
            \item 求 $\{a_n\}$ 的通项公式．
        \end{enumerate}

        \item （12分）

        某厂研制了一种生产高精产品的设备，为检验新设备生产产品的某项指标有无提高，用一台旧设备和一台新设备各生产了10件产品，得到数据．旧设备和新设备生产产品的该项指标的样本平均数分别记为 $\bar{x}$ 和 $\bar{y}$，样本方差分别记为 $s_1^2$ 和 $s_2^2$．

        \textsf{[插入表格：新旧设备数据]}

        \begin{enumerate}[label=(\arabic*)]
            \item 求 $\bar{x},\bar{y},s_1^2,s_2^2$；
            \item 判断新设备生产产品的该项指标的均值较旧设备是否有显著提高．
        \end{enumerate}

        \item （12分）

        已知直三棱柱 $ABC$-$A_1B_1C_1$ 中，侧面 $AA_1B_1B$ 为正方形，$AB=BC=2$，$E,F$ 分别为 $AC$ 和 $CC_1$ 的中点，$D$ 为棱 $A_1B_1$ 上的点，$BF\perp A_1B_1$．

        \textsf{[插入图片：直三棱柱示意图]}

        \begin{enumerate}[label=(\arabic*)]
            \item 证明：$BF\perp DE$；
            \item 当 $B_1D$ 为何值时，面 $BB_1C_1C$ 与面 $DFE$ 所成的二面角的正弦值最小？
        \end{enumerate}

        \item （12分）

        已知函数 $f(x)=(x-1)\e^x-ax^2+b$，曲线 $y=f(x)$ 在点 $(1,f(1))$ 处的切线方程为 $y=2\e x-\e$．

        \begin{enumerate}[label=(\arabic*)]
            \item 求 $a,b$ 的值；
            \item 讨论 $f(x)$ 的零点个数，并说明理由．
        \end{enumerate}

        \item （12分）

        抛物线 $C$ 的顶点为坐标原点 $O$，焦点在 $x$ 轴上，直线 $l:x=1$ 交 $C$ 于 $P,Q$ 两点，且 $OP\perp OQ$．已知点 $M(2,0)$，且 $\odot M$ 与 $l$ 相切．

        \begin{enumerate}[label=(\arabic*)]
            \item 求 $C$，$\odot M$ 的方程；
            \item 设 $A_1,A_2,A_3$ 是 $C$ 上的三个点，直线 $A_1A_2,A_1A_3$ 均与 $\odot M$ 相切．判断直线 $A_2A_3$ 与 $\odot M$ 的位置关系，并说明理由．
        \end{enumerate}

        \item （10分）[选修4-4：坐标系与参数方程]

        在直角坐标系 $xOy$ 中，$\odot C$ 的圆心为 $C(2,1)$，半径为 $1$．

        \begin{enumerate}[label=(\arabic*)]
            \item 写出 $\odot C$ 的一个参数方程；
            \item 过点 $F(4,1)$ 作 $\odot C$ 的两条切线，求这两条切线的极坐标方程．
        \end{enumerate}

        \item （10分）[选修4-5：不等式选讲]

        已知函数 $f(x)=|x-a|+|x+3|$．

        \begin{enumerate}[label=(\arabic*)]
            \item 当 $a=1$ 时，求不等式 $f(x)\geqslant6$ 的解集；
            \item 若 $f(x)>-a$，求 $a$ 的取值范围．
        \end{enumerate}
    \end{enumerate}
\end{description}
\clearpage

\subsection{2021年全国II卷（文）}\label{gk:2021年全国II卷（文）}

\begin{description}
    \item[一、选择题] 本题共12小题，每小题5分，共60分．在每小题给出的四个选项中，只有一项是符合题目要求的．
    \begin{enumerate}[leftmargin=5pt]
        \item 设集合 $M=\{1,2,3,4,5\}$，$N=\{2,4,6\}$，则 $M\cap N=$
       
        {\centering\begin{tabular*}{\linewidth}{@{}@{\extracolsep{\fill}}llll@{}}
          \(\text{A. }\{2,4\}\)  & \(\text{B. }\{1,3,5\}\)  & \(\text{C. }\{1,2,3,4,5,6\}\)  & \(\text{D. }\{6\}\)
        \end{tabular*}\par}

        \item 设 $\i z=4+3\i$，则 $z=$
       
        {\centering\begin{tabular*}{\linewidth}{@{}@{\extracolsep{\fill}}llll@{}}
          \(\text{A. }-3-4\i\)  & \(\text{B. }-3+4\i\)  & \(\text{C. }3-4\i\)  & \(\text{D. }3+4\i\)
        \end{tabular*}\par}

        \item 已知命题 $p:\exists x\in\mathbb{R},\sin x<1$；命题 $q:\forall x\in\mathbb{R},\e^{|x|}\geqslant1$，则下列命题中为真命题的是
       
        {\centering\begin{tabular*}{\linewidth}{@{}@{\extracolsep{\fill}}llll@{}}
          \(\text{A. }p\land q\)  & \(\text{B. }\neg p\land q\)  & \(\text{C. }p\land\neg q\)  & \(\text{D. }\neg(p\lor q)\)
        \end{tabular*}\par}

        \item 函数 $f(x)=\sin\dfrac{x}{3}+\cos\dfrac{x}{3}$ 的最小正周期和最大值分别是
       
        {\centering\begin{tabular*}{\linewidth}{@{}@{\extracolsep{\fill}}llll@{}}
          \(\text{A. }3\pi$和$\sqrt{2}\)  & \(\text{B. }3\pi$和$2\)  & \(\text{C. }6\pi$和$\sqrt{2}\)  & \(\text{D. }6\pi$和$2\)
        \end{tabular*}\par}

        \item 若 $x,y$ 满足约束条件 $\begin{cases}x+y\geqslant4,\\x-y\leqslant2,\\y\leqslant3,\end{cases}$ 则 $z=3x+y$ 的最小值为
       
        {\centering\begin{tabular*}{\linewidth}{@{}@{\extracolsep{\fill}}llll@{}}
          \(\text{A. }18\)  & \(\text{B. }10\)  & \(\text{C. }6\)  & \(\text{D. }4\)
        \end{tabular*}\par}

        \item $\cos^2\dfrac{\pi}{12}-\cos^2\dfrac{5\pi}{12}=$
       
        {\centering\begin{tabular*}{\linewidth}{@{}@{\extracolsep{\fill}}llll@{}}
          \(\text{A. }\dfrac{1}{2}\)  & \(\text{B. }\dfrac{\sqrt{3}}{3}\)  & \(\text{C. }\dfrac{\sqrt{2}}{2}\)  & \(\text{D. }\dfrac{\sqrt{3}}{2}\)
        \end{tabular*}\par}

        \item 在区间 $(0,\dfrac{1}{2})$ 随机取 $1$ 个数，则取到的数小于 $\dfrac{1}{3}$ 的概率为
       
        {\centering\begin{tabular*}{\linewidth}{@{}@{\extracolsep{\fill}}llll@{}}
          \(\text{A. }\dfrac{3}{4}\)  & \(\text{B. }\dfrac{2}{3}\)  & \(\text{C. }\dfrac{1}{3}\)  & \(\text{D. }\dfrac{1}{6}\)
        \end{tabular*}\par}

        \item 下列函数中最小值为 $4$ 的是
       
        {\centering\begin{tabular*}{\linewidth}{@{}@{\extracolsep{\fill}}llll@{}}
          \(\text{A. }y=x^2+2x+4\)  & \(\text{B. }y=|\sin x|+\dfrac{4}{|\sin x|}\)  & \(\text{C. }y=2^x+2^{2-x}\)  & \(\text{D. }y=\ln x+\dfrac{4}{\ln x}\)
        \end{tabular*}\par}

        \item 设函数 $f(x)=\dfrac{1-x}{1+x}$，则下列函数中为奇函数的是
       
        {\centering\begin{tabular*}{\linewidth}{@{}@{\extracolsep{\fill}}llll@{}}
          \(\text{A. }f(x-1)-1\)  & \(\text{B. }f(x-1)+1\)  & \(\text{C. }f(x+1)-1\)  & \(\text{D. }f(x+1)+1\)
        \end{tabular*}\par}

        \item 在正方体 $ABCD$-$A_1B_1C_1D_1$ 中，$P$ 为 $B_1D_1$ 的中点，则直线 $PB$ 与 $AD_1$ 所成的角为
       
        {\centering\begin{tabular*}{\linewidth}{@{}@{\extracolsep{\fill}}llll@{}}
          \(\text{A. }\dfrac{\pi}{2}\)  & \(\text{B. }\dfrac{\pi}{3}\)  & \(\text{C. }\dfrac{\pi}{4}\)  & \(\text{D. }\dfrac{\pi}{6}\)
        \end{tabular*}\par}

        \item 设 $B$ 是椭圆 $C:\dfrac{x^2}{5}+y^2=1$ 的上顶点，点 $P$ 在 $C$ 上，则 $|PB|$ 的最大值为
       
        {\centering\begin{tabular*}{\linewidth}{@{}@{\extracolsep{\fill}}llll@{}}
          \(\text{A. }\dfrac{5}{2}\)  & \(\text{B. }\sqrt{6}\)  & \(\text{C. }\sqrt{5}\)  & \(\text{D. }2\)
        \end{tabular*}\par}

        \item 设 $a\neq0$，若 $x=a$ 为函数 $f(x)=a(x-a)^2(x-b)$ 的极大值点，则
       
        {\centering\begin{tabular*}{\linewidth}{@{}@{\extracolsep{\fill}}llll@{}}
          \(\text{A. }a<b\)  & \(\text{B. }a>b\)  & \(\text{C. }ab<a^2\)  & \(\text{D. }ab>a^2\)
        \end{tabular*}\par}
    \end{enumerate}
\end{description}

\begin{description}
    \item[二、填空题] 本题共4小题，每小题5分，共20分．
    \begin{enumerate}[leftmargin=5pt]
      \addtocounter{enumi}{+12}
        \item 已知向量 $a=(2,5)$，$b=(\lambda,4)$，若 $a\parallel b$，则 $\lambda=$ \rule[-0.3ex]{3em}{0.4pt}．
        \item 双曲线 $\dfrac{x^2}{4}-\dfrac{y^2}{5}=1$ 的右焦点到直线 $x+2y-8=0$ 的距离为 \rule[-0.3ex]{3em}{0.4pt}．
        \item 记 $\triangle ABC$ 的内角 $A,B,C$ 的对边分别为 $a,b,c$，面积为 $\sqrt{3}$，$B=60^\circ$，$a^2+c^2=3ac$，则 $b=$ \rule[-0.3ex]{3em}{0.4pt}．
        \item 以图\ding{172}为正视图，在图\ding{173}\ding{174}\ding{175}\ding{176}中选两个分别作为侧视图和俯视图，组成某个三棱锥的三视图，则所选侧视图和俯视图的编号依次为 \rule[-0.3ex]{3em}{0.4pt}．
        
        \textsf{[插入图片：三视图选项]}
    \end{enumerate}
\end{description}

\begin{description}
    \item[三、解答题] 本题共7小题，共70分．解答应写出文字说明、证明过程或演算步骤．第17\textasciitilde21题为必考题，第22、23题为选考题．
    \begin{enumerate}[leftmargin=5pt]
      \addtocounter{enumi}{+16}
        \item （12分）

        记 $S_n$ 为数列 $\{a_n\}$ 的前 $n$ 项和，已知 $a_n>0$，$a_2=3a_1$，且数列 $\{\sqrt{S_n}\}$ 是等差数列，证明：$\{a_n\}$ 是等差数列．

        \item （12分）

        已知直三棱柱 $ABC$-$A_1B_1C_1$ 中，侧面 $AA_1B_1B$ 为正方形，$AB=BC=2$，$E,F$ 分别为 $AC$ 和 $CC_1$ 的中点，$BF\perp A_1B_1$．

        \textsf{[插入图片：直三棱柱示意图]}

        \begin{enumerate}[label=(\arabic*)]
            \item 求三棱锥 $F$-$EBC$ 的体积；
            \item 已知 $D$ 为棱 $A_1B_1$ 上的点，证明：$BF\perp DE$．
        \end{enumerate}

        \item （12分）

        设 $\{a_n\}$ 是首项为 $1$ 的等比数列，数列 $\{b_n\}$ 满足 $b_n=\dfrac{na_n}{3}$．已知 $a_1,3a_2,9a_3$ 成等差数列．

        \begin{enumerate}[label=(\arabic*)]
            \item 求 $\{a_n\}$ 和 $\{b_n\}$ 的通项公式；
            \item 记 $S_n$ 和 $T_n$ 分别为 $\{a_n\}$ 和 $\{b_n\}$ 的前 $n$ 项和．证明：$T_n<\dfrac{S_n}{2}$．
        \end{enumerate}

        \item （12分）

        已知抛物线 $C:y^2=2px\,(p>0)$ 的焦点 $F$ 到准线的距离为 $2$．

        \begin{enumerate}[label=(\arabic*)]
            \item 求 $C$ 的方程；
            \item 已知 $O$ 为坐标原点，点 $P$ 在 $C$ 上，点 $Q$ 满足 $\overrightarrow{PQ}=9\overrightarrow{QF}$，求直线 $OQ$ 斜率的最大值．
        \end{enumerate}

        \item （12分）

        已知函数 $f(x)=x^3-x^2+ax+1$．

        \begin{enumerate}[label=(\arabic*)]
            \item 讨论 $f(x)$ 的单调性；
            \item 求曲线 $y=f(x)$ 过坐标原点的切线与曲线 $y=f(x)$ 的公共点的坐标．
        \end{enumerate}

        \item （10分）[选修4-4：坐标系与参数方程]

        在直角坐标系 $xOy$ 中，$\odot C$ 的圆心为 $C(2,1)$，半径为 $1$．

        \begin{enumerate}[label=(\arabic*)]
            \item 写出 $\odot C$ 的一个参数方程；
            \item 过点 $F(4,1)$ 作 $\odot C$ 的两条切线，求这两条切线的极坐标方程．
        \end{enumerate}

        \item （10分）[选修4-5：不等式选讲]

        已知函数 $f(x)=|x-a|+|x+3|$．

        \begin{enumerate}[label=(\arabic*)]
            \item 当 $a=1$ 时，求不等式 $f(x)\geqslant6$ 的解集；
            \item 若 $f(x)>-a$，求 $a$ 的取值范围．
        \end{enumerate}
    \end{enumerate}
\end{description}
\clearpage

\subsection{2021年新高考I卷}\label{gk:2021年新高考I卷}

\begin{description}
    \item[一、单选题] 本题共8小题，每小题5分，共40分．在每小题给出的四个选项中，只有一项是符合题目要求的．
    \begin{enumerate}[leftmargin=5pt]
        \item 设集合 $A=\{x\mid -2<x<4\}$，$B=\{2,3,4,5\}$，则 $A\cap B=$ （ ）

        {\centering\begin{tabular*}{\linewidth}{@{}@{\extracolsep{\fill}}llll@{}}
          \(\text{A. }\{2\}\)  & \(\text{B. }\{2,3\}\)  & \(\text{C. }\{3,4\}\)  & \(\text{D. }\{2,3,4\}\)
        \end{tabular*}\par}

        \item 已知 $z=2-\i$，则 $z(\overline{z}+\i)=$ （ ）

        {\centering\begin{tabular*}{\linewidth}{@{}@{\extracolsep{\fill}}llll@{}}
          \(\text{A. }6-2\i\)  & \(\text{B. }4-2\i\)  & \(\text{C. }6+2\i\)  & \(\text{D. }4+2\i\)
        \end{tabular*}\par}

        \item 已知圆锥的底面半径为 $\sqrt{2}$，其侧面展开图为一个半圆，则该圆锥的母线长为 （ ）

        {\centering\begin{tabular*}{\linewidth}{@{}@{\extracolsep{\fill}}llll@{}}
          \(\text{A. }2\)  & \(\text{B. }2\sqrt{2}\)  & \(\text{C. }4\)  & \(\text{D. }4\sqrt{2}\)
        \end{tabular*}\par}

        \item 下列区间中，函数 $f(x)=7\sin\left(x-\frac{\pi}{6}\right)$ 单调递增的区间是 （ ）

        {\centering\begin{tabular*}{\linewidth}{@{}@{\extracolsep{\fill}}llll@{}}
          \(\text{A. }\left(0,\frac{\pi}{2}\right)\)  & \(\text{B. }\left(\frac{\pi}{2},\pi\right)\)  & \(\text{C. }\left(\pi,\frac{3\pi}{2}\right)\)  & \(\text{D. }\left(\frac{3\pi}{2},2\pi\right)\)
        \end{tabular*}\par}

        \item 已知 $F_{1}$，$F_{2}$ 是椭圆 $C:\frac{x^{2}}{9}+\frac{y^{2}}{4}=1$ 的两个焦点，点 $M$ 在 $C$ 上，则 $\left|MF_{1}\right|\cdot\left|MF_{2}\right|$ 的最大值为 （ ）

        {\centering\begin{tabular*}{\linewidth}{@{}@{\extracolsep{\fill}}llll@{}}
          \(\text{A. }13\)  & \(\text{B. }12\)  & \(\text{C. }9\)  & \(\text{D. }6\)
        \end{tabular*}\par}

        \item 若 $\tan\theta=-2$，则 $\frac{\sin\theta(1+\sin2\theta)}{\sin\theta+\cos\theta}=$ （ ）

        {\centering\begin{tabular*}{\linewidth}{@{}@{\extracolsep{\fill}}llll@{}}
          \(\text{A. }-\frac{6}{5}\)  & \(\text{B. }-\frac{2}{5}\)  & \(\text{C. }\frac{2}{5}\)  & \(\text{D. }\frac{6}{5}\)
        \end{tabular*}\par}

        \item 若过点 $(a,b)$ 可以作曲线 $y=\e^{x}$ 的两条切线，则 （ ）

        {\centering\begin{tabular*}{\linewidth}{@{}@{\extracolsep{\fill}}llll@{}}
          \(\text{A. }\e^{b}<a\)  & \(\text{B. }\e^{a}<b\)  & \(\text{C. }0<a<\e^{b}\)  & \(\text{D. }0<b<\e^{a}\)
        \end{tabular*}\par}

        \item 有6个相同的球，分别标有数字1,2,3,4,5,6，从中有放回的随机取两次，每次取1个球．甲表示事件"第一次取出的球的数字是1"，乙表示事件"第二次取出的球的数字是2"，丙表示事件"两次取出的球的数字之和是8"，丁表示事件"两次取出的球的数字之和是7"，则 （ ）

        {\centering\begin{tabular*}{\linewidth}{@{}@{\extracolsep{\fill}}llll@{}}
          \(\text{A. }\)甲与丙相互独立  & \(\text{B. }\)甲与丁相互独立  & \(\text{C. }\)乙与丙相互独立  & \(\text{D. }\)丙与丁相互独立
        \end{tabular*}\par}
    \end{enumerate}

    \item[二、多选题] 本题共4小题，每小题5分，共20分．在每小题给出的选项中，有多项符合题目要求．全部选对的得5分，部分选对的得2分，有选错的得0分．
    \begin{enumerate}[leftmargin=5pt]
      \addtocounter{enumi}{+8}
        \item 有一组样本数据 $x_{1},x_{2},\cdots,x_{n}$，由这组数据得到新样本数据 $y_{1},y_{2},\cdots,y_{n}$，其中 $y_{i}=x_{i}+c$（$i=1,2,\cdots,n$），$c$ 为非零常数，则 （ ）
        \begin{enumerate}[label=\Alph*.]
            \item 两组样本数据的样本平均数相同
            \item 两组样本数据的样本中位数相同
            \item 两组样本数据的样本标准差相同
            \item 两组样本数据的样本极差相同
        \end{enumerate}

        \item 已知 $O$ 为坐标原点，点 $P_{1}(\cos\alpha,\sin\alpha)$，$P_{2}(\cos\beta,-\sin\beta)$，$P_{3}(\cos(\alpha+\beta),\sin(\alpha+\beta))$，$A(1,0)$，则 （ ）
        \begin{enumerate}[label=\Alph*.]
            \item $\left|\overrightarrow{OP_{1}}\right|=\left|\overrightarrow{OP_{2}}\right|$
            \item $\left|\overrightarrow{AP_{1}}\right|=\left|\overrightarrow{AP_{2}}\right|$
            \item $\overrightarrow{OA}\cdot\overrightarrow{OP_{3}}=\overrightarrow{OP_{1}}\cdot\overrightarrow{OP_{2}}$
            \item $\overrightarrow{OA}\cdot\overrightarrow{OP_{1}}=\overrightarrow{OP_{2}}\cdot\overrightarrow{OP_{3}}$
        \end{enumerate}

        \item 已知点 $P$ 在圆 $(x-5)^{2}+(y-5)^{2}=16$ 上，点 $A(4,0)$，$B(0,2)$，则 （ ）
        \begin{enumerate}[label=\Alph*.]
            \item 点 $P$ 到直线 $AB$ 的距离小于10
            \item 点 $P$ 到直线 $AB$ 的距离大于2
            \item 当 $\angle PBA$ 最小时，$|PB|=3\sqrt{2}$
            \item 当 $\angle PBA$ 最大时，$|PB|=3\sqrt{2}$
        \end{enumerate}

        \item 在正三棱柱 $ABC$-$A_{1}B_{1}C_{1}$ 中，$AB=AA_{1}=1$，点 $P$ 满足 $\overrightarrow{BP}=\lambda\overrightarrow{BC}+\mu\overrightarrow{BB_{1}}$，其中 $\lambda\in[0,1]$，$\mu\in[0,1]$，则 （ ）
        \begin{enumerate}[label=\Alph*.]
            \item 当 $\lambda=1$ 时，$\triangle AB_{1}P$ 的周长为定值
            \item 当 $\mu=1$ 时，三棱锥 $P$-$A_{1}BC$ 的体积为定值
            \item 当 $\lambda=\frac{1}{2}$ 时，有且仅有一个点 $P$，使得 $A_{1}P\perp BP$
            \item 当 $\mu=\frac{1}{2}$ 时，有且仅有一个点 $P$，使得 $A_{1}B\perp$ 平面 $AB_{1}P$
        \end{enumerate}
    \end{enumerate}

    \item[三、填空题] 本题共4小题，每小题5分，共20分．
    \begin{enumerate}[leftmargin=5pt]
      \addtocounter{enumi}{+12}
        \item 已知函数 $f(x)=x^{3}(a\cdot2^{x}-2^{-x})$ 是偶函数，则 $a=$ \rule[-0.3ex]{3em}{0.4pt}．
        \item 已知 $O$ 为坐标原点，抛物线 $C:y^{2}=2px$（$p>0$）的焦点为 $F$，$P$ 为 $C$ 上一点，$PF$ 与 $x$ 轴垂直，$Q$ 为 $x$ 轴上一点，且 $PQ\perp OP$．若 $|FQ|=6$，则 $C$ 的准线方程为 \rule[-0.3ex]{3em}{0.4pt}．
        \item 函数 $f(x)=|2x-1|-2\ln x$ 的最小值为 \rule[-0.3ex]{3em}{0.4pt}．
        \item 某校学生在研究民间剪纸艺术时，发现剪纸时经常会沿纸的某条对称轴把纸对折．规格为 $20\,\text{dm}\times 12\,\text{dm}$ 的长方形纸，对折1次共可以得到 $10\,\text{dm}\times 12\,\text{dm}$，$20\,\text{dm}\times 6\,\text{dm}$ 两种规格的图形，它们的面积之和 $S_{1}=240\,\text{dm}^{2}$，对折2次共可以得到 $5\,\text{dm}\times 12\,\text{dm}$，$10\,\text{dm}\times 6\,\text{dm}$，$20\,\text{dm}\times 3\,\text{dm}$ 三种规格的图形，它们的面积之和 $S_{2}=180\,\text{dm}^{2}$，以此类推．则对折4次共可以得到不同规格图形的种数为 \rule[-0.3ex]{3em}{0.4pt}；如果对折 $n$ 次，那么 $\displaystyle\sum_{k=1}^{n}S_{k}=$ \rule[-0.3ex]{3em}{0.4pt} $\text{dm}^{2}$．
    \end{enumerate}

    \item[四、解答题] 本题共6小题，共70分．解答应写出文字说明、证明过程或演算步骤．
    \begin{enumerate}[leftmargin=5pt]
      \addtocounter{enumi}{+16}
        \item （10分）

        已知数列 $\{a_{n}\}$ 满足 $a_{1}=1$，$a_{n+1}=\begin{cases}a_{n}+1,&n\text{为奇数},\\ a_{n}+2,&n\text{为偶数}.\end{cases}$

        \begin{enumerate}[label=(\arabic*)]
            \item 记 $b_{n}=a_{2n}$，写出 $b_{1},b_{2}$，并求数列 $\{b_{n}\}$ 的通项公式；
            \item 求 $\{a_{n}\}$ 的前20项和．
        \end{enumerate}

        \item （12分）

        某学校组织"一带一路"知识竞赛，有A，B两类问题．每位参加比赛的同学先在两类问题中选择一类并从中随机抽取一个问题回答，若回答错误则该同学比赛结束；若回答正确则从另一类问题中再随机抽取一个问题回答，无论回答正确与否，该同学比赛结束．A类问题中的每个问题回答正确得20分，否则得0分；B类问题中的每个问题回答正确得80分，否则得0分．已知小明能正确回答A类问题的概率为0.8，能正确回答B类问题的概率为0.6，且能正确回答问题的概率与回答次序无关．

        \begin{enumerate}[label=(\arabic*)]
            \item 若小明先回答A类问题，记 $X$ 为小明的累计得分，求 $X$ 的分布列；
            \item 为使累计得分的期望最大，小明应选择先回答哪类问题？并说明理由．
        \end{enumerate}

        \item （12分）

        记 $\triangle ABC$ 的内角 $A,B,C$ 的对边分别为 $a,b,c$．已知 $b^{2}=ac$，点 $D$ 在边 $AC$ 上，$BD\sin\angle ABC=a\sin C$．

        \begin{enumerate}[label=(\arabic*)]
            \item 证明：$BD=b$；
            \item 若 $AD=2DC$，求 $\cos\angle ABC$．
        \end{enumerate}

        \item （12分）

        如图，在三棱锥 $A$-$BCD$ 中，平面 $ABD\perp$ 平面 $BCD$，$AB=AD$，$O$ 为 $BD$ 的中点．

        \begin{enumerate}[label=(\arabic*)]
            \item 证明：$OA\perp CD$；
            \item 若 $\triangle OCD$ 是边长为1的等边三角形，点 $E$ 在棱 $AD$ 上，$DE=2EA$，且二面角 $E$-$BC$-$D$ 的大小为 $45^{\circ}$，求三棱锥 $A$-$BCD$ 的体积．
        \end{enumerate}

        \item （12分）

        在平面直角坐标系 $xOy$ 中，已知点 $F_{1}(-\sqrt{17},0)$，$F_{2}(\sqrt{17},0)$，点 $M$ 满足 $\left|MF_{1}\right|-\left|MF_{2}\right|=2$．记 $M$ 的轨迹为 $C$．

        \textsf{[插入图片：双曲线示意图]}

        \begin{enumerate}[label=(\arabic*)]
            \item 求 $C$ 的方程；
            \item 设点 $T$ 在直线 $x=\frac{1}{2}$ 上，过 $T$ 的两条直线分别交 $C$ 于 $A,B$ 两点和 $P,Q$ 两点，且 $|TA|\cdot|TB|=|TP|\cdot|TQ|$，求直线 $AB$ 的斜率与直线 $PQ$ 的斜率之和．
        \end{enumerate}

        \item （12分）

        已知函数 $f(x)=x(1-\ln x)$．

        \begin{enumerate}[label=(\arabic*)]
            \item 讨论 $f(x)$ 的单调性；
            \item 设 $a,b$ 为两个不相等的正数，且 $b\ln a-a\ln b=a-b$，证明：$2<\frac{1}{a}+\frac{1}{b}<\e$．
        \end{enumerate}
    \end{enumerate}
\end{description}
\clearpage

\subsection{2021年新高考II卷}\label{gk:2021年新高考II卷}

\begin{description}
    \item[一、单选题] 本题共8小题，每小题5分，共40分．在每小题给出的四个选项中，只有一项是符合题目要求的．
    \begin{enumerate}[leftmargin=5pt]
        \item 复数 $\frac{2-\i}{1-3\i}$ 在复平面内对应的点所在的象限为 （ ）

        {\centering\begin{tabular*}{\linewidth}{@{}@{\extracolsep{\fill}}llll@{}}
          \(\text{A. }\)第一象限  & \(\text{B. }\)第二象限  & \(\text{C. }\)第三象限  & \(\text{D. }\)第四象限
        \end{tabular*}\par}

        \item 若全集 $U=\{1,2,3,4,5,6\}$，集合 $A=\{1,3,6\}$，$B=\{2,3,4\}$，则 $A\cap(\complement_{U}B)=$ （ ）

        {\centering\begin{tabular*}{\linewidth}{@{}@{\extracolsep{\fill}}llll@{}}
          \(\text{A. }\{3\}\)  & \(\text{B. }\{1,6\}\)  & \(\text{C. }\{5,6\}\)  & \(\text{D. }\{1,3\}\)
        \end{tabular*}\par}

        \item 若抛物线 $y^{2}=2px$（$p>0$）的焦点到直线 $y=x+1$ 的距离为 $\sqrt{2}$，则 $p=$ （ ）

        {\centering\begin{tabular*}{\linewidth}{@{}@{\extracolsep{\fill}}llll@{}}
          \(\text{A. }1\)  & \(\text{B. }2\)  & \(\text{C. }2\sqrt{2}\)  & \(\text{D. }4\)
        \end{tabular*}\par}

        \item 北斗三号全球卫星导航系统是我国航天事业的重要成果．在卫星导航系统中，地球静止同步轨道卫星的轨道位于地球赤道所在平面，轨道高度为 $36000\,\text{km}$（轨道高度是指卫星到地球表面的最短距离），把地球看成一个球心为 $O$，半径 $r$ 为 $6400\,\text{km}$ 的球，其上点 $A$ 的纬度是指 $OA$ 与赤道所在平面所成角的度数．地球表面能直接观测到的一颗地球静止同步轨道卫星的点的纬度的最大值记为 $\alpha$，该卫星信号覆盖的地球表面面积 $S=2\pi r^{2}(1-\cos\alpha)$（单位：$\text{km}^{2}$），则 $S$ 占地球表面积的百分比约为 （ ）

        {\centering\begin{tabular*}{\linewidth}{@{}@{\extracolsep{\fill}}llll@{}}
          \(\text{A. }26\%\)  & \(\text{B. }34\%\)  & \(\text{C. }42\%\)  & \(\text{D. }50\%\)
        \end{tabular*}\par}

        \item 正四棱台的上、下底面的边长为2,4，侧棱长为2，则其体积为 （ ）

        {\centering\begin{tabular*}{\linewidth}{@{}@{\extracolsep{\fill}}llll@{}}
          \(\text{A. }56\)  & \(\text{B. }28\sqrt{2}\)  & \(\text{C. }\frac{56}{3}\)  & \(\text{D. }\frac{28\sqrt{2}}{3}\)
        \end{tabular*}\par}

        \item 某物理量的测量结果服从正态分布 $N(10,\sigma^{2})$，则下列结论中不正确的是 （ ）

        {\centering\begin{tabular*}{\linewidth}{@{}@{\extracolsep{\fill}}llll@{}}
          \(\text{A. }\)$\sigma$ 越小，该物理量一次测量结果落在 $(9.9,10.1)$ 内的概率越大  & \(\text{B. }\)该物理量一次测量结果大于10的概率为0.5  & \(\text{C. }\)该物理量一次测量结果小于9.99的概率与大于10.01的概率相等  & \(\text{D. }\)该物理量一次测量结果落在 $(9.9,10.2)$ 内的概率与落在 $(10,10.3)$ 内的概率相等
        \end{tabular*}\par}

        \item 若 $a=\log_{5}2$，$b=\log_{8}3$，$c=\frac{1}{2}$，则 （ ）

        {\centering\begin{tabular*}{\linewidth}{@{}@{\extracolsep{\fill}}llll@{}}
          \(\text{A. }c<b<a\)  & \(\text{B. }b<a<c\)  & \(\text{C. }a<c<b\)  & \(\text{D. }a<b<c\)
        \end{tabular*}\par}

        \item 设函数 $f(x)$ 的定义域为 $\mathbb{R}$，且 $f(x+2)$ 为偶函数，$f(2x+1)$ 为奇函数，则 （ ）

        {\centering\begin{tabular*}{\linewidth}{@{}@{\extracolsep{\fill}}llll@{}}
          \(\text{A. }f\left(-\frac{1}{2}\right)=0\)  & \(\text{B. }f(-1)=0\)  & \(\text{C. }f(2)=0\)  & \(\text{D. }f(4)=0\)
        \end{tabular*}\par}
    \end{enumerate}

    \item[二、多选题] 本题共4小题，每小题5分，共20分．在每小题给出的选项中，有多项符合题目要求．全部选对的得5分，部分选对的得2分，有选错的得0分．
    \begin{enumerate}[leftmargin=5pt]
      \addtocounter{enumi}{+8}
        \item 下列统计量中可用于度量样本 $x_{1},x_{2},\cdots,x_{n}$ 离散程度的有 （ ）
        \begin{enumerate}[label=\Alph*.]
            \item $x_{1},x_{2},\cdots,x_{n}$ 的标准差
            \item $x_{1},x_{2},\cdots,x_{n}$ 的中位数
            \item $x_{1},x_{2},\cdots,x_{n}$ 的极差
            \item $x_{1},x_{2},\cdots,x_{n}$ 的平均数
        \end{enumerate}

        \item 如图，下列各正方体中，$O$ 为下底面的中心，$M,N$ 为顶点，$P$ 为所在棱的中点，则满足 $MN\perp OP$ 的是 （ ）

        \textsf{[插入图片：四个正方体示意图]}

        \begin{enumerate}[label=\Alph*.]
            \item 图A
            \item 图B
            \item 图C
            \item 图D
        \end{enumerate}

        \item 已知直线 $l:ax+by-r^{2}=0$ 与圆 $C:x^{2}+y^{2}=r^{2}$，点 $A(a,b)$，则下列说法正确的是 （ ）
        \begin{enumerate}[label=\Alph*.]
            \item 若点 $A$ 在圆 $C$ 上，则直线 $l$ 与圆 $C$ 相切
            \item 若点 $A$ 在圆 $C$ 内，则直线 $l$ 与圆 $C$ 相离
            \item 若点 $A$ 在圆 $C$ 外，则直线 $l$ 与圆 $C$ 相离
            \item 若点 $A$ 在直线 $l$ 上，则直线 $l$ 与圆 $C$ 相切
        \end{enumerate}

        \item 设正整数 $n=a_{0}\cdot2^{0}+a_{1}\cdot2^{1}+\cdots+a_{k-1}\cdot2^{k-1}+a_{k}\cdot2^{k}$，其中 $a_{i}\in\{0,1\}$，记 $\omega(n)=a_{0}+a_{1}+\cdots+a_{k}$．则 （ ）
        \begin{enumerate}[label=\Alph*.]
            \item $\omega(2n)=\omega(n)$
            \item $\omega(2n+3)=\omega(n)+1$
            \item $\omega(8n+5)=\omega(4n+3)$
            \item $\omega(2^{n}-1)=n$
        \end{enumerate}
    \end{enumerate}

    \item[三、填空题] 本题共4小题，每小题5分，共20分．
    \begin{enumerate}[leftmargin=5pt]
      \addtocounter{enumi}{+12}
        \item 已知双曲线 $C:\frac{x^{2}}{a^{2}}-\frac{y^{2}}{b^{2}}=1$（$a>0,b>0$）的离心率 $e=2$，则双曲线 $C$ 的渐近线方程为 \rule[-0.3ex]{3em}{0.4pt}．
        \item 写出一个同时具有下列性质\textcircled{1}\textcircled{2}\textcircled{3}的函数 $f(x)=$ \rule[-0.3ex]{3em}{0.4pt}．

        \textcircled{1} $f(x_{1}x_{2})=f(x_{1})f(x_{2})$；

        \textcircled{2} 当 $x\in(0,+\infty)$ 时，$f'(x)>0$；

        \textcircled{3} $f'(x)$ 是奇函数．
        \item 已知向量 $\bm{a}+\bm{b}+\bm{c}=\bm{0}$，$|\bm{a}|=1$，$|\bm{b}|=|\bm{c}|=2$，则 $\bm{a}\cdot\bm{b}+\bm{b}\cdot\bm{c}+\bm{c}\cdot\bm{a}=$ \rule[-0.3ex]{3em}{0.4pt}．
    \end{enumerate}

    \item[四、解答题] 本题共6小题，共70分．解答应写出文字说明、证明过程或演算步骤．
    \begin{enumerate}[leftmargin=5pt]
      \addtocounter{enumi}{+15}
        \item 已知函数 $f(x)=|\e^{x}-1|$，$x_{1}<0$，$x_{2}>0$，函数 $f(x)$ 的图象在点 $A(x_{1},f(x_{1}))$ 和点 $B(x_{2},f(x_{2}))$ 处的两条切线互相垂直，且分别交 $y$ 轴于 $M,N$ 两点，则 $\frac{|AM|}{|BN|}$ 的取值范围是 \rule[-0.3ex]{3em}{0.4pt}．

        \item （10分）

        记 $S_{n}$ 是公差不为0的等差数列 $\{a_{n}\}$ 的前 $n$ 项和，若 $a_{3}=S_{5}$，$a_{2}a_{4}=S_{4}$．

        \begin{enumerate}[label=(\arabic*)]
            \item 求数列 $\{a_{n}\}$ 的通项公式；
            \item 求使 $S_{n}>a_{n}$ 成立的 $n$ 的最小值．
        \end{enumerate}

        \item （12分）

        在 $\triangle ABC$ 中，角 $A,B,C$ 所对的边分别为 $a,b,c$，$b=a+1$，$c=a+2$．

        \begin{enumerate}[label=(\arabic*)]
            \item 若 $2\sin C=3\sin A$，求 $\triangle ABC$ 的面积；
            \item 是否存在正整数 $a$，使得 $\triangle ABC$ 为钝角三角形？若存在，求 $a$；若不存在，说明理由．
        \end{enumerate}

        \item （12分）

        如图，在四棱锥 $Q$-$ABCD$ 中，底面 $ABCD$ 是正方形，$AD=2$，$QD=QA=\sqrt{5}$，$QC=3$．

        \begin{enumerate}[label=(\arabic*)]
            \item 证明：平面 $QAD\perp$ 平面 $ABCD$；
            \item 求二面角 $B$-$QD$-$A$ 的余弦值．
        \end{enumerate}

        \item （12分）

        已知椭圆 $C:\frac{x^{2}}{a^{2}}+\frac{y^{2}}{b^{2}}=1$（$a>b>0$），右焦点为 $F(\sqrt{2},0)$，其离心率为 $\frac{\sqrt{6}}{3}$．

        \begin{enumerate}[label=(\arabic*)]
            \item 求椭圆 $C$ 的方程；
            \item 设 $M,N$ 是椭圆 $C$ 上的两点，直线 $MN$ 与曲线 $x^{2}+y^{2}=b^{2}$（$x>0$）相切．证明：$M,N,F$ 三点共线的充要条件是 $|MN|=\sqrt{3}$．
        \end{enumerate}

        \item （12分）

        一种微生物群体可以经过自身繁殖不断生存下来，设一个这种微生物为第0代，经过一次繁殖后为第1代，再经过一次繁殖后为第2代……该微生物每代繁殖的个数是相互独立的且有相同的分布列，设 $X$ 表示1个微生物个体繁殖下一代的个数，$P(X=i)=p_{i}>0$（$i=0,1,2,3$）．

        \begin{enumerate}[label=(\arabic*)]
            \item 已知 $p_{0}=0.4$，$p_{1}=0.3$，$p_{2}=0.2$，$p_{3}=0.1$，求 $E(X)$；
            \item 设 $p$ 表示该种微生物经过多代繁殖后临近灭绝的概率，$p$ 是关于 $x$ 的方程 $p_{0}+p_{1}x+p_{2}x^{2}+p_{3}x^{3}=x$ 的一个最小正实根，求证：当 $E(X)\leqslant1$ 时 $p=1$，当 $E(X)>1$ 时 $p<1$；
            \item 根据你的理解说明（2）问结论的实际含义．
        \end{enumerate}

        \item （12分）

        已知函数 $f(x)=(x-1)\e^{x}-ax^{2}+b$．

        \begin{enumerate}[label=(\arabic*)]
            \item 讨论 $f(x)$ 的单调性；
            \item 从下面两个条件中任选一个作为已知条件，证明：$f(x)$ 有一个零点．

            \ding{172} $\frac{1}{2}<a\leqslant\frac{\e^{2}}{2}$，$b>2a$；

            \ding{173} $0<a<\frac{1}{2}$，$b\leqslant2a$．
        \end{enumerate}
    \end{enumerate}
\end{description}
\clearpage


\subsection{2021年北京卷}\label{gk:2021年北京卷}

\begin{description}
    \item[一、选择题] 本题共10小题，每小题4分，共40分．在每小题给出的四个选项中，选出符合题目要求的一项．
    \begin{enumerate}[leftmargin=5pt]
        \item 已知集合 $A=\{x\mid -1<x<1\}$，$B=\{x\mid0\leqslant x\leqslant2\}$，则 $A\cup B=$
       
        {\centering\begin{tabular*}{\linewidth}{@{}@{\extracolsep{\fill}}llll@{}}
          \(\text{A. }\{x\mid -1<x<2\}\)  & \(\text{B. }\{x\mid -1<x\leqslant2\}\)  & \(\text{C. }\{x\mid0\leqslant x<1\}\)  & \(\text{D. }\{x\mid0\leqslant x\leqslant2\}\)
        \end{tabular*}\par}

        \item 在复平面内，复数 $z$ 满足 $(1-\i)z=2$，则 $z=$
       
        
        {\centering\begin{tabular*}{\linewidth}{@{}@{\extracolsep{\fill}}llll@{}}
          \(\text{A. }1+\i\)  & \(\text{B. }1-\i\)  & \(\text{C. }-1+\i\)  & \(\text{D. }-1-\i\)
        \end{tabular*}\par}

        \item 设函数 $f(x)$ 的定义域为 $[0,1]$，则"$f(x)$ 在 $[0,1]$ 上单调递增"是"$f(x)$ 在 $[0,1]$ 上的最大值为 $f(1)$"的
       
        {\centering\begin{tabular*}{\linewidth}{@{}@{\extracolsep{\fill}}llll@{}}
          \(\text{A. }\)充分不必要条件  & \(\text{B. }\)必要不充分条件  & \(\text{C. }\)充要条件  & \(\text{D. }\)既不充分也不必要条件
        \end{tabular*}\par}

        \item 某四面体的三视图如图所示，该四面体的表面积为
        
        \textsf{[插入图片：三视图]}

        {\centering\begin{tabular*}{\linewidth}{@{}@{\extracolsep{\fill}}llll@{}}
          \(\text{A. }\dfrac{3+\sqrt{3}}{2}\)  & \(\text{B. }3+\sqrt{3}\)  & \(\text{C. }\dfrac{1}{2}+\sqrt{3}\)  & \(\text{D. }\dfrac{3}{2}+\sqrt{3}\)
        \end{tabular*}\par}

        \item 双曲线 $C:\dfrac{x^2}{a^2}-\dfrac{y^2}{b^2}=1$ 过点 $(\sqrt{2},\sqrt{3})$，且离心率为 $2$，则该双曲线的标准方程为
       
        {\centering\begin{tabular*}{\linewidth}{@{}@{\extracolsep{\fill}}llll@{}}
          \(\text{A. }x^2-\dfrac{y^2}{3}=1\)  & \(\text{B. }\dfrac{x^2}{3}-y^2=1\)  & \(\text{C. }x^2-\dfrac{\sqrt{3}y^2}{3}=1\)  & \(\text{D. }\dfrac{x^2}{3}-\dfrac{y^2}{3}=1\)
        \end{tabular*}\par}

        \item 已知 $\{a_n\}$ 和 $\{b_n\}$ 是两个等差数列，且 $\dfrac{a_n}{b_n}=\dfrac{3n-1}{2n+3}\,(n\in\mathbb{N}^*)$，则 $\dfrac{a_7}{b_7}=$
       
        {\centering\begin{tabular*}{\linewidth}{@{}@{\extracolsep{\fill}}llll@{}}
          \(\text{A. }\dfrac{20}{17}\)  & \(\text{B. }\dfrac{13}{10}\)  & \(\text{C. }\dfrac{37}{17}\)  & \(\text{D. }\dfrac{40}{17}\)
        \end{tabular*}\par}

        \item 已知函数 $f(x)=\cos x-\cos2x$，则该函数
       
        {\centering\begin{tabular*}{\linewidth}{@{}@{\extracolsep{\fill}}llll@{}}
          \(\text{A. }\)是奇函数，最大值为 $2$  & \(\text{B. }\)是偶函数，最大值为 $2$ \\
          \(\text{C. }\)是奇函数，最大值为 $\dfrac{9}{8}$  & \(\text{D. }\)是偶函数，最大值为 $\dfrac{9}{8}$
        \end{tabular*}\par}

        \item 对24小时内降水在平地上的积水厚度（mm）进行如下定义：$0\sim10$为小雨，$10\sim25$为中雨，$25\sim50$为大雨，$50\sim100$为暴雨．小明用一个圆锥形容器接了24小时的雨水，则这一天的雨水属于哪个等级？
        
        \textsf{[插入图片：圆锥形容器示意图]}

        {\centering\begin{tabular*}{\linewidth}{@{}@{\extracolsep{\fill}}llll@{}}
          \(\text{A. }\)小雨  & \(\text{B. }\)中雨  & \(\text{C. }\)大雨  & \(\text{D. }\)暴雨
        \end{tabular*}\par}

        \item 已知直线 $y=kx+m\,(m\neq0)$ 与圆 $(x-5)^2+y^2=9$ 和圆 $x^2+y^2=16$ 均相切，则 $k=$
       
        {\centering\begin{tabular*}{\linewidth}{@{}@{\extracolsep{\fill}}llll@{}}
          \(\text{A. }\pm\dfrac{3}{4}\)  & \(\text{B. }\pm\dfrac{4}{3}\)  & \(\text{C. }\pm\dfrac{3}{4}$或$\pm\dfrac{4}{3}\)  & \(\text{D. }\pm\dfrac{3}{4}$或$\pm\dfrac{7}{24}\)
        \end{tabular*}\par}

        \item 已知 $a_n=2n-1$，若集合 $M=\{k\mid a_k a_{k+1}\geqslant a_{k+2},1\leqslant k\leqslant20\}$ 中恰有 $n$ 个元素，则 $n=$
       
        {\centering\begin{tabular*}{\linewidth}{@{}@{\extracolsep{\fill}}llll@{}}
          \(\text{A. }10\)  & \(\text{B. }11\)  & \(\text{C. }12\)  & \(\text{D. }13\)
        \end{tabular*}\par}
    \end{enumerate}
\end{description}

\begin{description}
    \item[二、填空题] 本题共5小题，每小题5分，共25分．
    \begin{enumerate}[leftmargin=5pt]
      \addtocounter{enumi}{+10}
        \item 在 $(x^3-\dfrac{1}{x})^4$ 的展开式中，常数项为 \rule[-0.3ex]{3em}{0.4pt}．
        \item 已知抛物线 $C:y^2=4x$，$C$ 的焦点为 $F$，点 $M$ 在 $C$ 上，且 $|FM|=6$，则 $M$ 的横坐标是 \rule[-0.3ex]{3em}{0.4pt}．
        \item 已知 $a=(2,1)$，$b=(2,-1)$，$c=(0,1)$，则 $(a+b)\cdot c=$ \rule[-0.3ex]{3em}{0.4pt}；$a\cdot b=$ \rule[-0.3ex]{3em}{0.4pt}．
        \item 若 $P(\cos\theta,\sin\theta)$ 与 $Q(\cos(\theta+\dfrac{\pi}{6}),\sin(\theta+\dfrac{\pi}{6}))$ 关于 $y$ 轴对称，写出一个符合题意的 $\theta=$ \rule[-0.3ex]{3em}{0.4pt}．
        \item 已知函数 $f(x)=|\lg x|-kx-2$，给出下列四个结论：

        \ding{172}当 $k=0$ 时，$f(x)$ 恰有 $2$ 个零点；
        \ding{173}存在 $k\in(0,1)$，使得 $f(x)$ 恰有 $1$ 个零点；
        \ding{174}存在 $k\in(0,1)$，使得 $f(x)$ 恰有 $3$ 个零点；
        \ding{175}对任意 $k<0$，$f(x)$ 恰有 $1$ 个零点．

        其中所有正确结论的序号是 \rule[-0.3ex]{3em}{0.4pt}．
    \end{enumerate}
\end{description}

\begin{description}
    \item[三、解答题] 本题共6小题，共85分．解答应写出文字说明、证明过程或演算步骤．
    \begin{enumerate}[leftmargin=5pt]
      \addtocounter{enumi}{+15}
        \item （13分）

        在 $\triangle ABC$ 中，$c=2b\cos B$，$C=\dfrac{2\pi}{3}$．

        \begin{enumerate}[label=(\arabic*)]
            \item 求 $\angle B$；
            \item 在条件\ding{172}、条件\ding{173}、条件\ding{174}中任选一个作为已知，使 $\triangle ABC$ 存在且唯一确定，并求 $BC$ 边上的中线长．
            \begin{quote}
            条件\ding{172}：$c=\sqrt{3}b$；条件\ding{173}：$\triangle ABC$ 的周长为 $4+2\sqrt{3}$；条件\ding{174}：$\triangle ABC$ 的面积为 $\dfrac{3\sqrt{3}}{4}$．
            \end{quote}
        \end{enumerate}

        \item （14分）

        已知正方体 $ABCD$-$A_1B_1C_1D_1$，点 $E$ 为 $A_1D_1$ 的中点，直线 $B_1C_1$ 交平面 $CDE$ 于点 $F$．

        \textsf{[插入图片：正方体示意图]}

        \begin{enumerate}[label=(\arabic*)]
            \item 求证：点 $F$ 为 $B_1C_1$ 的中点；
            \item 若点 $M$ 为棱 $A_1B_1$ 上一点，且二面角 $M$-$CF$-$E$ 的余弦值为 $\dfrac{\sqrt{5}}{3}$，求 $\dfrac{A_1M}{A_1B_1}$ 的值．
        \end{enumerate}

        \item （14分）

        在核酸检测中，"$k$合1"混采核酸检测是指：先将 $k$ 个人的样本混合在一起进行1次检测，如果这 $k$ 个人都没感染新冠病毒，则检测结果为阴性；如果这 $k$ 个人中有人感染新冠病毒，则检测结果阳性，此时需对每人再进行一次检测，得到每人的检测结果，检测结束．现对100人进行核酸检测，假设其中只有2人感染新冠病毒，并假设每次检测结果准确．

        \begin{enumerate}[label=(\arabic*)]
            \item 将这100人随机分成10组，每组10人，且对每组都采用"10合1"混采核酸检测．
            \begin{enumerate}[label=(\roman*)]
                \item 如果感染新冠病毒的2人在同一组，求检测的总次数；
                \item 已知感染新冠病毒的2人分在同一组的概率为 $\dfrac{1}{11}$．设 $X$ 是检测的总次数，求 $X$ 的分布列与数学期望 $E(X)$．
            \end{enumerate}
            \item 将这100人随机分成20组，每组5人，且对每组都采用"5合1"混采核酸检测．设 $Y$ 是检测的总次数，试判断数学期望 $E(Y)$ 与(1)中 $E(X)$ 的大小．（结论不要求证明）
        \end{enumerate}

        \item （15分）

        已知函数 $f(x)=\dfrac{3-2x}{x^2+a}$．

        \begin{enumerate}[label=(\arabic*)]
            \item 若 $a=0$，求 $y=f(x)$ 在 $(1,f(1))$ 处的切线方程；
            \item 若函数 $f(x)$ 在 $x=-1$ 处取得极值，求 $f(x)$ 的单调区间，以及最大值和最小值．
        \end{enumerate}

        \item （15分）

        已知椭圆 $E:\dfrac{x^2}{a^2}+\dfrac{y^2}{b^2}=1\,(a>b>0)$ 过点 $A(0,-2)$，以四个顶点围成的四边形面积为 $4\sqrt{5}$．

        \begin{enumerate}[label=(\arabic*)]
            \item 求椭圆 $E$ 的标准方程；
            \item 过点 $P(0,-3)$ 的直线 $l$ 斜率为 $k$，交椭圆 $E$ 于不同的两点 $B,C$，直线 $AB,AC$ 交 $y=-3$ 于点 $M,N$，若 $|PM|+|PN|\leqslant15$，求 $k$ 的取值范围．
        \end{enumerate}

        \item （15分）

        设 $p$ 为实数．若无穷数列 $\{a_n\}$ 满足如下三个性质，则称 $\{a_n\}$ 为 $\Re_p$ 数列：
        \begin{quote}
        \ding{172}$a_1+p\geqslant0$，且 $a_2+p=0$；
        \ding{173}$a_{4n-1}<a_{4n}\,(n=1,2,\cdots)$；
        \ding{174}$a_{m+n}\in\{a_m+a_n+p,a_m+a_n+p+1\}\,(m,n=1,2,\cdots)$．
        \end{quote}

        \begin{enumerate}[label=(\arabic*)]
            \item 如果 $\{a_n\}$ 的前4项为 $2,-2,-2,-1$，那么 $\{a_n\}$ 是否可能为 $\Re_2$ 数列？说明理由；
            \item 若 $\{a_n\}$ 是 $\Re_0$ 数列，求 $a_5$；
            \item 设 $\{a_n\}$ 的各项均为正数，且 $\{a_n\}$ 是 $\Re_p$ 数列，求证：对任意 $n\in\mathbb{N}^*$，存在 $k\in\mathbb{N}^*$，使得 $a_{n+k}>a_n$．
        \end{enumerate}
    \end{enumerate}
\end{description}
\clearpage








\subsection{2021年上海卷}\label{gk:2021年上海卷}

\begin{description}
    \item[一、填空题] 本大题共有12题，满分54分，第1\textasciitilde6题每题4分，第7\textasciitilde12题每题5分．
    \begin{enumerate}[leftmargin=5pt]
        \item 已知 $z_1=1+\i$，$z_2=2+3\i$，则 $z_1+z_2=$ \rule[-0.3ex]{3em}{0.4pt}．
        \item 已知 $A=\{x\mid2x\leqslant1\}$，$B=\{-1,0,1\}$，则 $A\cap B=$ \rule[-0.3ex]{3em}{0.4pt}．
        \item 若 $x^2+y^2-2x-4y=0$，则圆心坐标为 \rule[-0.3ex]{3em}{0.4pt}．
        \item 如图，正方形 $ABCD$ 的边长为 $3$，则 $\overrightarrow{AB}\cdot\overrightarrow{AC}=$ \rule[-0.3ex]{3em}{0.4pt}．
        \item 已知 $f(x)=\dfrac{3}{x}+2$，则 $f^{-1}(1)=$ \rule[-0.3ex]{3em}{0.4pt}．
        \item 已知 $(x+a)^5$ 的展开式中，$x^2$ 的系数为 $80$，则 $a=$ \rule[-0.3ex]{3em}{0.4pt}．
        \item 已知 $\begin{cases}x=3t-4,\\y=t+2\end{cases}$ 则 $x-y=$ \rule[-0.3ex]{3em}{0.4pt}．
        \item 已知 $\{a_n\}$ 为无穷等比数列，$a_1=3$，$a_n$ 的各项和为 $9$，$b_n=a_{2n}$，则数列 $\{b_n\}$ 的各项和为 \rule[-0.3ex]{3em}{0.4pt}．
        \item 已知圆柱的底面半径为 $1$，高为 $2$，$AB$ 为上底面圆的直径，$C$ 为下底面圆周上的一个动点，则 $\triangle ABC$ 的面积的取值范围为 \rule[-0.3ex]{3em}{0.4pt}．
        \item 已知花博会有四个不同的场馆 $A,B,C,D$，甲、乙两人每人选 $2$ 个去参观，则两人选择中恰有一个场馆相同的选法有 \rule[-0.3ex]{3em}{0.4pt} 种．
        \item 已知抛物线 $y^2=2px\,(p>0)$，若第一象限的 $A,B$ 在抛物线上，焦点为 $F$，$|AF|=2$，$|BF|=4$，$|AB|=3$，则直线 $AB$ 的斜率为 \rule[-0.3ex]{3em}{0.4pt}．
        \item 已知 $a_i\in\mathbb{N}^*\,(i=1,2,\cdots,9)$，对任意的 $k\in\mathbb{N}^*\,(2\leqslant k\leqslant8)$，$a_k=a_{k-1}+1$ 或 $a_k=a_{k+1}-1$ 中有且仅有一个成立，且 $a_1=6$，$a_9=9$，则 $a_1+a_2+\cdots+a_9$ 的最小值为 \rule[-0.3ex]{3em}{0.4pt}．
    \end{enumerate}
\end{description}

\begin{description}
    \item[二、选择题] 本题共4小题，每小题5分，共20分．
    \begin{enumerate}[leftmargin=5pt]
      \addtocounter{enumi}{+12}
        \item 以下哪个函数既是奇函数，又是减函数
       
        {\centering\begin{tabular*}{\linewidth}{@{}@{\extracolsep{\fill}}llll@{}}
          \(\text{A. }y=-x^3\)  & \(\text{B. }y=x^3\)  & \(\text{C. }y=\log_3 x\)  & \(\text{D. }y=3^x\)
        \end{tabular*}\par}

        \item 已知参数方程 $\begin{cases}x=3t^2-1,\\y=2t\end{cases}$，$t\in[-1,1]$，下列选项的图中，符合该方程的曲线是
        
        \textsf{[插入图片：曲线四选一]}

        \item 已知 $f(x)=3\sin x+2$，存在 $x_1,x_2$，对任意 $x\in\mathbb{R}$，$f(x_1)\leqslant f(x)\leqslant f(x_2)$，则 $|x_1-x_2|$ 的最小值为
       
        {\centering\begin{tabular*}{\linewidth}{@{}@{\extracolsep{\fill}}llll@{}}
          \(\text{A. }\dfrac{\pi}{2}\)  & \(\text{B. }\pi\)  & \(\text{C. }2\pi\)  & \(\text{D. }4\pi\)
        \end{tabular*}\par}

        \item 已知 $x_1,y_1,x_2,y_2,x_3,y_3$ 同时满足\ding{172}$x_1<y_1,x_2<y_2,x_3<y_3$；\ding{173}$x_1+y_1=x_2+y_2=x_3+y_3$；\ding{174}$x_1y_1+x_3y_3=2x_2y_2$，以下哪个选项恒成立
       
        {\centering\begin{tabular*}{\linewidth}{@{}@{\extracolsep{\fill}}llll@{}}
          \(\text{A. }2x_2<x_1+x_3\)  & \(\text{B. }2x_2>x_1+x_3\)  & \(\text{C. }x_2^2<x_1x_3\)  & \(\text{D. }x_2^2>x_1x_3\)
        \end{tabular*}\par}
    \end{enumerate}
\end{description}

\begin{description}
    \item[三、解答题] 本题共5小题，共76分．
    \begin{enumerate}[leftmargin=5pt]
      \addtocounter{enumi}{+16}
        \item （14分）

        如图，在长方体 $ABCD$-$A_1B_1C_1D_1$ 中，已知 $AB=BC=2$，$AA_1=3$．

        \textsf{[插入图片：长方体示意图]}

        \begin{enumerate}[label=(\arabic*)]
            \item 若 $P$ 是棱 $A_1D_1$ 上的动点，求三棱锥 $C$-$PAD$ 的体积；
            \item 求直线 $AB_1$ 与平面 $ACC_1A_1$ 所成的角的大小．
        \end{enumerate}

        \item （14分）

        在 $\triangle ABC$ 中，已知 $a=3$，$b=2c$．

        \begin{enumerate}[label=(\arabic*)]
            \item 若 $A=\dfrac{2\pi}{3}$，求 $S_{\triangle ABC}$；
            \item 若 $2\sin B-\sin C=1$，求 $C_{\triangle ABC}$．
        \end{enumerate}

        \item （14分）

        已知某企业今年（2021年）第一个季度的营业额为 $1.1$ 亿元，以后每个季度的营业额比上个季度增加 $0.05$ 亿元，该企业第一季度的利润为 $0.16$ 亿元，以后每季度比前一季度增长 $4\%$．

        \begin{enumerate}[label=(\arabic*)]
            \item 求2021年起前20个季度营业额的总和；
            \item 请问哪一季度的利润第一次超过该季度营业额的 $18\%$？
        \end{enumerate}

        \item （17分）

        已知椭圆 $\Gamma:\dfrac{x^2}{2}+y^2=1$，$F_1,F_2$ 是 $\Gamma$ 的左、右焦点，$P$ 是 $\Gamma$ 上的点（不与顶点重合）．

        \begin{enumerate}[label=(\arabic*)]
            \item 求 $\Gamma$ 的离心率；
            \item 过 $F_2$ 作 $\angle F_1PF_2$ 的外角平分线的垂线，求垂足 $Q$ 的轨迹方程；
            \item 设 $A$ 为椭圆的右顶点，是否存在 $y$ 轴上的定点 $M$，使得以 $PA$ 为直径的圆恒过定点 $M$？若存在，求出 $M$，若不存在，说明理由．
        \end{enumerate}

        \item （17分）

        已知 $f(x)$ 是定义在 $\mathbb{R}$ 上的函数，若对任意 $x_1,x_2\in\mathbb{R}$，$x_1-x_2\in S$，都有 $f(x_1)-f(x_2)\in S$，则称 $f(x)$ 是 $S$ 关联的．

        \begin{enumerate}[label=(\arabic*)]
            \item 判断 $f(x)=2x-1$ 是否是 $[1,+\infty)$ 关联的；
            \item 若 $f(x)$ 是 $\{3\}$ 关联的，且 $x\in[0,3)$ 时，$f(x)=x^2$，求不等式 $2<f(x)<3$ 的解集；
            \item 证明："$f(x)$ 是 $\{1\}$ 关联的，且是 $[0,+\infty)$ 关联的"当且仅当"$f(x)$ 是 $[1,+\infty)$ 关联的"．
        \end{enumerate}
    \end{enumerate}
\end{description}
\clearpage

\subsection{2021年天津卷}\label{gk:2021年天津卷}

\begin{description}
    \item[一、选择题] 本题共9小题，每小题5分，共45分．在每小题给出的四个选项中，只有一项是符合题目要求的．
    \begin{enumerate}[leftmargin=5pt]
        \item 设集合 $A=\{-1,0,1\}$，$B=\{1,3,5\}$，$C=\{0,2,4\}$，则 $(A\cap B)\cup C=$
       
        {\centering\begin{tabular*}{\linewidth}{@{}@{\extracolsep{\fill}}llll@{}}
          \(\text{A. }\{0\}\)  & \(\text{B. }\{0,1,3,5\}\)  & \(\text{C. }\{0,1,2,4\}\)  & \(\text{D. }\{0,2,3,4\}\)
        \end{tabular*}\par}

        \item 已知 $a\in\mathbb{R}$，则"$a>6$"是"$a^2>36$"的
       
        {\centering\begin{tabular*}{\linewidth}{@{}@{\extracolsep{\fill}}llll@{}}
          \(\text{A. }\)充分不必要条件  & \(\text{B. }\)必要不充分条件  & \(\text{C. }\)充要条件  & \(\text{D. }\)既不充分也不必要条件
        \end{tabular*}\par}

        \item 函数 $y=\dfrac{\ln|x|}{x^2+2}$ 的图像大致为
        
        \textsf{[插入图片：函数图象四选一]}

        \item 从某网络平台推荐的影视作品中抽取400部，统计其评分数据，将所得400个评分数据分为8组：$[66,70),[70,74),\cdots,[94,98]$，并整理得到如下的频率分布直方图，则评分在区间 $[82,86)$ 内的影视作品数量是
        
        \textsf{[插入图片：频率分布直方图]}

        {\centering\begin{tabular*}{\linewidth}{@{}@{\extracolsep{\fill}}llll@{}}
          \(\text{A. }20\)  & \(\text{B. }40\)  & \(\text{C. }64\)  & \(\text{D. }80\)
        \end{tabular*}\par}

        \item 设 $a=\log_2 0.3$，$b=\log_{\frac{1}{2}}0.4$，$c=0.4^{0.3}$，则
       
        {\centering\begin{tabular*}{\linewidth}{@{}@{\extracolsep{\fill}}llll@{}}
          \(\text{A. }a<b<c\)  & \(\text{B. }c<a<b\)  & \(\text{C. }b<c<a\)  & \(\text{D. }a<c<b\)
        \end{tabular*}\par}

        \item 两个圆锥的底面是一个球的同一截面上的两个圆，顶点均在球面上，若球的体积为 $\dfrac{32\pi}{3}$，两个圆锥的高之比为 $1:3$，则这两个圆锥的体积之和为
       
        {\centering\begin{tabular*}{\linewidth}{@{}@{\extracolsep{\fill}}llll@{}}
          \(\text{A. }3\pi\)  & \(\text{B. }4\pi\)  & \(\text{C. }9\pi\)  & \(\text{D. }12\pi\)
        \end{tabular*}\par}

        \item 若 $2^a=5^b=10$，则 $\dfrac{1}{a}+\dfrac{1}{b}=$
       
        {\centering\begin{tabular*}{\linewidth}{@{}@{\extracolsep{\fill}}llll@{}}
          \(\text{A. }-1\)  & \(\text{B. }\ln7\)  & \(\text{C. }1\)  & \(\text{D. }\log_7 10\)
        \end{tabular*}\par}

        \item 已知双曲线 $\dfrac{x^2}{a^2}-\dfrac{y^2}{b^2}=1\,(a>0,b>0)$ 的右焦点与抛物线 $y^2=2px\,(p>0)$ 的焦点重合，抛物线的准线交双曲线于 $A,B$ 两点，交双曲线的渐近线于 $C,D$ 两点，若 $|CD|=\sqrt{2}|AB|$，则双曲线的离心率为
       
        {\centering\begin{tabular*}{\linewidth}{@{}@{\extracolsep{\fill}}llll@{}}
          \(\text{A. }\sqrt{2}\)  & \(\text{B. }\sqrt{3}\)  & \(\text{C. }2\)  & \(\text{D. }3\)
        \end{tabular*}\par}

        \item 设 $a\in\mathbb{R}$，函数 $f(x)=\begin{cases}\cos(2\pi x-2\pi a),&x<a,\\x^2-2(a+1)x+a^2+5,&x\geqslant a.\end{cases}$ 若 $f(x)$ 在区间 $(0,+\infty)$ 内恰有 $6$ 个零点，则 $a$ 的取值范围是
       
        {\centering\begin{tabular*}{\linewidth}{@{}@{\extracolsep{\fill}}llll@{}}
          \(\text{A. }\bigl(2,\dfrac{9}{4}\bigr]\cup\bigl(\dfrac{5}{2},\dfrac{11}{4}\bigr]\)  & \(\text{B. }\bigl(\dfrac{7}{4},2\bigr]\cup\bigl(\dfrac{5}{2},\dfrac{11}{4}\bigr]\) \\
          \(\text{C. }\bigl(2,\dfrac{9}{4}\bigr]\cup\bigl[\dfrac{11}{4},3\bigr)\)  & \(\text{D. }\bigl(\dfrac{7}{4},2\bigr]\cup\bigl[\dfrac{11}{4},3\bigr)\)
        \end{tabular*}\par}
    \end{enumerate}
\end{description}

\begin{description}
    \item[二、填空题] 本题共6小题，每小题5分，共30分．
    \begin{enumerate}[leftmargin=5pt]
      \addtocounter{enumi}{+9}
        \item $\i$ 是虚数单位，复数 $\dfrac{9+2\i}{2+\i}=$ \rule[-0.3ex]{3em}{0.4pt}．
        \item 在 $\bigl(2x^3+\dfrac{1}{x}\bigr)^6$ 的展开式中，$x^6$ 的系数是 \rule[-0.3ex]{3em}{0.4pt}．
        \item 若斜率为 $\sqrt{3}$ 的直线与 $y$ 轴交于点 $A$，与圆 $x^2+(y-1)^2=1$ 相切于点 $B$，则 $|AB|=$ \rule[-0.3ex]{3em}{0.4pt}．
        \item 若 $a>0,b>0$，则 $\dfrac{1}{a}+\dfrac{a}{b^2}+b$ 的最小值为 \rule[-0.3ex]{3em}{0.4pt}．
        \item 甲、乙两人在每次猜谜活动中各猜一个谜语，若一方猜对且另一方猜错，则猜对的一方获胜，否则本次平局．已知每次活动中，甲、乙猜对的概率分别为 $\dfrac{5}{6}$ 和 $\dfrac{1}{5}$，且每次活动中甲、乙猜对与否互不影响，各次活动也互不影响，则一次活动中，甲获胜的概率是 \rule[-0.3ex]{3em}{0.4pt}；3次活动中，甲至少获胜2次的概率为 \rule[-0.3ex]{3em}{0.4pt}．
        \item 在边长为 $1$ 的等边三角形 $ABC$ 中，$D$ 为线段 $BC$ 上的动点，$DE\perp AB$ 且交 $AB$ 于点 $E$，$DF\parallel AB$ 且交 $AC$ 于点 $F$，则 $|2\overrightarrow{BE}+\overrightarrow{DF}|$ 的值为 \rule[-0.3ex]{3em}{0.4pt}；$(\overrightarrow{DE}+\overrightarrow{DF})\cdot\overrightarrow{DA}$ 的最小值为 \rule[-0.3ex]{3em}{0.4pt}．
    \end{enumerate}
\end{description}

\begin{description}
    \item[三、解答题] 本题共5小题，共75分．解答应写出文字说明、证明过程或演算步骤．
    \begin{enumerate}[leftmargin=5pt]
      \addtocounter{enumi}{+15}
        \item （14分）

        在 $\triangle ABC$ 中，角 $A,B,C$ 所对的边分别为 $a,b,c$．已知 $\sin A:\sin B:\sin C=2:1:\sqrt{2}$，$b=\sqrt{2}$．

        \begin{enumerate}[label=(\arabic*)]
            \item 求 $a$ 的值；
            \item 求 $\cos C$ 的值；
            \item 求 $\sin\bigl(2C-\dfrac{\pi}{6}\bigr)$ 的值．
        \end{enumerate}

        \item （15分）

        如图，在棱长为 $2$ 的正方体 $ABCD$-$A_1B_1C_1D_1$ 中，$E$ 为棱 $BC$ 的中点，$F$ 为棱 $CD$ 的中点．

        \textsf{[插入图片：正方体示意图]}

        \begin{enumerate}[label=(\arabic*)]
            \item 求证：$D_1F\parallel$ 平面 $A_1EC_1$；
            \item 求直线 $AC_1$ 与平面 $A_1EC_1$ 所成角的正弦值；
            \item 求二面角 $A$-$A_1C_1$-$E$ 的正弦值．
        \end{enumerate}

        \item （15分）

        已知椭圆 $\dfrac{x^2}{a^2}+\dfrac{y^2}{b^2}=1\,(a>b>0)$ 的右焦点为 $F$，上顶点为 $B$，离心率为 $\dfrac{2\sqrt{5}}{5}$，且 $|BF|=\sqrt{5}$．

        \begin{enumerate}[label=(\arabic*)]
            \item 求椭圆的方程；
            \item 直线 $l$ 与椭圆有唯一的公共点 $M$，与 $y$ 轴的正半轴交于点 $N$，过 $N$ 与 $BF$ 垂直的直线交 $x$ 轴于点 $P$．若 $MP\parallel BF$，求直线 $l$ 的方程．
        \end{enumerate}

        \item （15分）

        已知 $\{a_n\}$ 是公差为 $2$ 的等差数列，其前8项的和为 $64$．$\{b_n\}$ 是公比大于 $0$ 的等比数列，$b_1=4$，$b_3-b_2=48$．

        \begin{enumerate}[label=(\arabic*)]
            \item 求 $\{a_n\}$ 和 $\{b_n\}$ 的通项公式；
            \item 记 $c_n=b_{2n}+\dfrac{1}{b_n}$，$n\in\mathbb{N}^*$．
            \begin{enumerate}[label=(\roman*)]
                \item 证明：$\{c_n^2-c_{2n}\}$ 是等比数列；
                \item 求 $\displaystyle\sum_{k=1}^{n}\sqrt{\dfrac{a_k a_{k+1}}{c_k^2-c_{2k}}}$．
            \end{enumerate}
        \end{enumerate}

        \item （16分）

        已知 $a>0$，函数 $f(x)=ax-x\e^x$．

        \begin{enumerate}[label=(\arabic*)]
            \item 求曲线 $y=f(x)$ 在点 $(0,f(0))$ 处的切线方程；
            \item 证明 $f(x)$ 存在唯一的极值点；
            \item 若存在 $a$，使得 $f(x)\leqslant a+b$ 对任意 $x\in\mathbb{R}$ 成立，求实数 $b$ 的取值范围．
        \end{enumerate}
    \end{enumerate}
\end{description}
\clearpage


\subsection{2021年浙江卷}\label{gk:2021年浙江卷}

\begin{description}
    \item[一、单选题] 本题共10小题，每小题4分，共40分．在每小题给出的四个选项中，只有一项是符合题目要求的．
    \begin{enumerate}[leftmargin=5pt]
        \item 设集合 $A=\{x\mid x\geqslant1\}$，$B=\{x\mid-1<x<2\}$，则 $A\cap B=$ （ ）

        {\centering\begin{tabular*}{\linewidth}{@{}@{\extracolsep{\fill}}llll@{}}
          \(\text{A. }\{x\mid x>-1\}\)  & \(\text{B. }\{x\mid x\geqslant1\}\)  & \(\text{C. }\{x\mid-1<x<1\}\)  & \(\text{D. }\{x\mid1\leqslant x<2\}\)
        \end{tabular*}\par}

        \item 已知 $a\in\mathbb{R}$，$(1+a\i)\i=3+\i$（$\i$ 为虚数单位），则 $a=$ （ ）

        {\centering\begin{tabular*}{\linewidth}{@{}@{\extracolsep{\fill}}llll@{}}
          \(\text{A. }-1\)  & \(\text{B. }1\)  & \(\text{C. }-3\)  & \(\text{D. }3\)
        \end{tabular*}\par}

        \item 已知非零向量 $\bm{a},\bm{b},\bm{c}$，则"$\bm{a}\cdot\bm{c}=\bm{b}\cdot\bm{c}$"是"$\bm{a}=\bm{b}$"的 （ ）

        {\centering\begin{tabular*}{\linewidth}{@{}@{\extracolsep{\fill}}llll@{}}
          \(\text{A. }\)充分不必要条件  & \(\text{B. }\)必要不充分条件  & \(\text{C. }\)充要条件  & \(\text{D. }\)既不充分也不必要条件
        \end{tabular*}\par}

        \item 某几何体的三视图如图所示（单位：cm），则该几何体的体积（单位：$\text{cm}^{3}$）是 （ ）

        \textsf{[插入图片：正视图、侧视图、俯视图]}

        {\centering\begin{tabular*}{\linewidth}{@{}@{\extracolsep{\fill}}llll@{}}
          \(\text{A. }\frac{3}{2}\)  & \(\text{B. }3\)  & \(\text{C. }\frac{3\sqrt{2}}{2}\)  & \(\text{D. }3\sqrt{2}\)
        \end{tabular*}\par}

        \item 若实数 $x,y$ 满足约束条件 $\begin{cases}x+1\geqslant0,\\ x-y\leqslant0,\\ 2x+3y-1\leqslant0,\end{cases}$ 则 $z=x-\frac{1}{2}y$ 的最小值是 （ ）

        {\centering\begin{tabular*}{\linewidth}{@{}@{\extracolsep{\fill}}llll@{}}
          \(\text{A. }-2\)  & \(\text{B. }-\frac{3}{2}\)  & \(\text{C. }-\frac{1}{2}\)  & \(\text{D. }\frac{1}{10}\)
        \end{tabular*}\par}

        \item 如图，已知正方体 $ABCD$-$A_{1}B_{1}C_{1}D_{1}$，$M,N$ 分别是 $A_{1}D$，$D_{1}B$ 的中点，则 （ ）

        \textsf{[插入图片：正方体示意图]}

        {\centering\begin{tabular*}{\linewidth}{@{}@{\extracolsep{\fill}}llll@{}}
          \(\text{A. }\)直线 $A_{1}D$ 与直线 $D_{1}B$ 垂直，直线 $MN\parallel$ 平面 $ABCD$  & \(\text{B. }\)直线 $A_{1}D$ 与直线 $D_{1}B$ 平行，直线 $MN\perp$ 平面 $BDD_{1}B_{1}$  & \(\text{C. }\)直线 $A_{1}D$ 与直线 $D_{1}B$ 相交，直线 $MN\parallel$ 平面 $ABCD$  & \(\text{D. }\)直线 $A_{1}D$ 与直线 $D_{1}B$ 异面，直线 $MN\perp$ 平面 $BDD_{1}B_{1}$
        \end{tabular*}\par}

        \item 已知函数 $f(x)=x^{2}+\frac{1}{4}$，$g(x)=\sin x$，则图象为如图的函数可能是 （ ）

        \textsf{[插入图片：函数图象]}

        {\centering\begin{tabular*}{\linewidth}{@{}@{\extracolsep{\fill}}llll@{}}
          \(\text{A. }y=f(x)+g(x)-\frac{1}{4}\)  & \(\text{B. }y=f(x)-g(x)-\frac{1}{4}\)  & \(\text{C. }y=f(x)g(x)\)  & \(\text{D. }y=\frac{g(x)}{f(x)}\)
        \end{tabular*}\par}

        \item 已知 $\alpha,\beta,\gamma$ 是互不相同的锐角，则在 $\sin\alpha\cos\beta$，$\sin\beta\cos\gamma$，$\sin\gamma\cos\alpha$ 三个值中，大于 $\frac{1}{2}$ 的个数的最大值是 （ ）

        {\centering\begin{tabular*}{\linewidth}{@{}@{\extracolsep{\fill}}llll@{}}
          \(\text{A. }0\)  & \(\text{B. }1\)  & \(\text{C. }2\)  & \(\text{D. }3\)
        \end{tabular*}\par}

        \item 已知 $a,b\in\mathbb{R}$，$ab>0$，函数 $f(x)=ax^{2}+b$（$x\in\mathbb{R}$）．若 $f(s-t)$，$f(s)$，$f(s+t)$ 成等比数列，则平面上点 $(s,t)$ 的轨迹是 （ ）

        {\centering\begin{tabular*}{\linewidth}{@{}@{\extracolsep{\fill}}llll@{}}
          \(\text{A. }\)直线和圆  & \(\text{B. }\)直线和椭圆  & \(\text{C. }\)直线和双曲线  & \(\text{D. }\)直线和抛物线
        \end{tabular*}\par}

        \item 已知数列 $\{a_{n}\}$ 满足 $a_{1}=1$，$a_{n+1}=\frac{a_{n}}{1+\sqrt{a_{n}}}$（$n\in\mathbb{N}^{*}$）．记数列 $\{a_{n}\}$ 的前 $n$ 项和为 $S_{n}$，则 （ ）

        {\centering\begin{tabular*}{\linewidth}{@{}@{\extracolsep{\fill}}llll@{}}
          \(\text{A. }\frac{3}{2}<S_{100}<3\)  & \(\text{B. }3<S_{100}<4\)  & \(\text{C. }4<S_{100}<\frac{9}{2}\)  & \(\text{D. }\frac{9}{2}<S_{100}<5\)
        \end{tabular*}\par}
    \end{enumerate}

    \item[二、填空题] 本题共7小题，共36分．
    \begin{enumerate}[leftmargin=5pt]
      \addtocounter{enumi}{+10}
        \item 我国古代数学家赵爽用弦图给出了勾股定理的证明．弦图是由四个全等的直角三角形和中间的一个小正方形拼成的一个大正方形（如图所示）．若直角三角形直角边的长分别是3,4，记大正方形的面积为 $S_{1}$，小正方形的面积为 $S_{2}$，则 $\frac{S_{1}}{S_{2}}=$ \rule[-0.3ex]{3em}{0.4pt}．

        \textsf{[插入图片：弦图]}

        \item 已知 $a\in\mathbb{R}$，函数 $f(x)=\begin{cases}x^{2}-4,&x>2,\\|x-3|+a,&x\leqslant2.\end{cases}$ 若 $f(f(\sqrt{6}))=3$，则 $a=$ \rule[-0.3ex]{3em}{0.4pt}．
        \item 已知多项式 $(x-1)^{3}+(x+1)^{4}=x^{4}+a_{1}x^{3}+a_{2}x^{2}+a_{3}x+a_{4}$，则 $a_{1}=$ \rule[-0.3ex]{3em}{0.4pt}；$a_{2}+a_{3}+a_{4}=$ \rule[-0.3ex]{3em}{0.4pt}．
        \item 在 $\triangle ABC$ 中，$\angle B=60^{\circ}$，$AB=2$，$M$ 是 $BC$ 的中点，$AM=2\sqrt{3}$，则 $AC=$ \rule[-0.3ex]{3em}{0.4pt}；$\cos\angle MAC=$ \rule[-0.3ex]{3em}{0.4pt}．
        \item 袋中有4个红球，$m$ 个黄球，$n$ 个绿球．现从中任取两个球，记取出的红球数为 $\xi$，若取出的两个球都是红球的概率为 $\frac{1}{6}$，一红一黄的概率为 $\frac{1}{3}$，则 $m-n=$ \rule[-0.3ex]{3em}{0.4pt}，$E(\xi)=$ \rule[-0.3ex]{3em}{0.4pt}．
        \item 已知椭圆 $\frac{x^{2}}{a^{2}}+\frac{y^{2}}{b^{2}}=1$（$a>b>0$），焦点 $F_{1}(-c,0)$，$F_{2}(c,0)$（$c>0$），若过 $F_{1}$ 的直线和圆 $\left(x-\frac{1}{2}c\right)^{2}+y^{2}=c^{2}$ 相切，与椭圆在第一象限交于点 $P$，且 $PF_{2}\perp x$ 轴，则该直线的斜率是 \rule[-0.3ex]{3em}{0.4pt}，椭圆的离心率是 \rule[-0.3ex]{3em}{0.4pt}．
        \item 已知平面向量 $\bm{a},\bm{b},\bm{c}$（$\bm{c}\neq\bm{0}$）满足 $|\bm{a}|=1$，$|\bm{b}|=2$，$\bm{a}\cdot\bm{b}=0$，$(\bm{a}-\bm{b})\cdot\bm{c}=0$．记向量 $\bm{d}$ 在 $\bm{a},\bm{b}$ 方向上的投影分别为 $x,y$，$\bm{d}-\bm{a}$ 在 $\bm{c}$ 方向上的投影为 $z$，则 $x^{2}+y^{2}+z^{2}$ 的最小值是 \rule[-0.3ex]{3em}{0.4pt}．
    \end{enumerate}

    \item[三、解答题] 本题共5小题，共74分．解答应写出文字说明、证明过程或演算步骤．
    \begin{enumerate}[leftmargin=5pt]
      \addtocounter{enumi}{+17}
        \item （14分）

        设函数 $f(x)=\sin x+\cos x$（$x\in\mathbf{R}$）．

        \begin{enumerate}[label=(\arabic*)]
            \item 求函数 $y=\left[f\left(x+\frac{\pi}{2}\right)\right]^{2}$ 的最小正周期；
            \item 求函数 $y=f(x)f\left(x-\frac{\pi}{4}\right)$ 在 $\left[0,\frac{\pi}{2}\right]$ 上的最大值．
        \end{enumerate}

        \item （15分）

        如图，在四棱锥 $P$-$ABCD$ 中，底面 $ABCD$ 是平行四边形，$\angle ABC=120^{\circ}$，$AB=1$，$BC=4$，$PA=\sqrt{15}$，$M,N$ 分别为 $BC,PC$ 的中点，$PD\perp DC$，$PM\perp MD$．

        \textsf{[插入图片：四棱锥示意图]}

        \begin{enumerate}[label=(\arabic*)]
            \item 证明：$AB\perp PM$；
            \item 求直线 $AN$ 与平面 $PDM$ 所成角的正弦值．
        \end{enumerate}

        \item （15分）

        已知数列 $\{a_{n}\}$ 的前 $n$ 项和为 $S_{n}$，$a_{1}=-\frac{9}{4}$，且 $4S_{n+1}=3S_{n}-9$（$n\in\mathbb{N}^{*}$）．

        \begin{enumerate}[label=(\arabic*)]
            \item 求数列 $\{a_{n}\}$ 的通项公式；
            \item 设数列 $\{b_{n}\}$ 满足 $3b_{n}+(n-4)a_{n}=0$，记 $\{b_{n}\}$ 的前 $n$ 项和为 $T_{n}$，若 $T_{n}\leqslant\lambda b_{n}$ 对任意 $n\in\mathbb{N}^{*}$ 恒成立，求实数 $\lambda$ 的取值范围．
        \end{enumerate}

        \item （15分）

        如图，已知 $F$ 是抛物线 $y^{2}=2px$（$p>0$）的焦点，$M$ 是抛物线的准线与 $x$ 轴的交点，且 $|MF|=2$．

        \textsf{[插入图片：抛物线示意图]}

        \begin{enumerate}[label=(\arabic*)]
            \item 求抛物线方程；
            \item 设过点 $F$ 的直线交抛物线于 $A$、$B$ 两点，斜率为2的直线 $l$ 与直线 $MA,MB,AB,x$ 轴依次交于点 $P,Q,R,N$，且 $|RN|^{2}=|PN|\cdot|QN|$，求直线 $l$ 在 $x$ 轴上截距的范围．
        \end{enumerate}

        \item （15分）

        设 $a,b$ 为实数，且 $a>1$，函数 $f(x)=a^{x}-bx+\e^{2}$（$x\in\mathbb{R}$）．

        \begin{enumerate}[label=(\arabic*)]
            \item 求函数 $f(x)$ 的单调区间；
            \item 若对任意 $b>2\e^{2}$，函数 $f(x)$ 有两个不同的零点，求 $a$ 的取值范围；
            \item 当 $a=\e$ 时，证明：对任意 $b>\e^{4}$，函数 $f(x)$ 有两个不同的零点 $x_{1},x_{2}$，满足 $x_{2}>\frac{b\ln b}{2\e^{2}}x_{1}+\frac{\e^{2}}{b}$．
        \end{enumerate}
    \end{enumerate}
\end{description}
\clearpage
